\documentclass[11pt,a4paper]{article}
\usepackage{amsmath,amssymb,amsthm,bm,booktabs,geometry,hyperref,mathtools,xcolor}
\usepackage{array,longtable}
\usepackage{kotex}
\geometry{margin=25mm}
\emergencystretch=2em
\hypersetup{colorlinks=true,linkcolor=blue,citecolor=blue,urlcolor=blue}

\newcommand{\Tr}{\operatorname{Tr}}
\newcommand{\dd}{\mathrm{d}}
\newcommand{\vk}{\mathbf{k}}
\newcommand{\vR}{\mathbf{R}}
\newcommand{\vq}{\mathbf{q}}
\newcommand{\vG}{\mathbf{G}}
\newcommand{\vT}{\mathbf{T}}
\newcommand{\vc}{\mathbf{c}}
\newcommand{\gradR}[2]{\frac{\partial #1}{\partial R_{#2}}}
\newcommand{\Herm}{\mathrm{H}}

\theoremstyle{definition}
\newtheorem{definition}{정의}[section]
\newtheorem{remark}{주의}[section]

\title{Periodic Density-Corrected DFT Analytic Gradients}
\author{PBC DC-DFT Package Project}
\date{Active scope: FFTDF/GDF/RSDF, target multigrid, PBE/SCAN/r2SCAN}

\begin{document}
\maketitle
\begin{abstract}
이 문서는 restricted closed-shell periodic Hartree--Fock density에서 semilocal target
functional을 평가하는 density-corrected DFT의 energy와 fixed-cell nuclear analytic
gradient를 유도한다. 핵심은 target energy가 HF orbital에 대해 stationary하지 않다는
점이며, 이로 인해 HF CPHF orbital Hessian과 adjoint Z-vector가 필요하다. FFTDF,
Gaussian density fitting(GDF), range-separated density fitting(RSDF), GTH
pseudopotential, Gamma 및 sampled k-point, PBE와 SCAN/r2SCAN의 수식과 구현 convention을
한 흐름으로 정리한다. Source GDF/RSDF와 독립적으로 선택하는 target multigrid,
matrix-free MINRES Z-vector 및 bounded-memory response algebra도 포함한다. Stress, cell
derivative, multigrid PP와 open-shell response는 이 판의 active 범위가 아니다.
\end{abstract}
\tableofcontents

\section{목적, 전제와 읽는 방법}
\label{sec:scope}

\subsection{계산하려는 물리량}

원자핵 좌표를 모은 벡터를 $\vR$라 하자. 이 자료의 중심 대상은 periodic
restricted Hartree--Fock(HF) density matrix $D_{\mathrm{HF}}(\vR)$를 먼저 얻은 뒤,
같은 density에서 별도의 semilocal density functional approximation(DFA)을 평가한
energy
\begin{equation}
  E_{\mathrm{DC-DFA}}(\vR)
  = E_{\mathrm{DFA}}\!\left[D_{\mathrm{HF}}(\vR);\vR\right]
  \label{eq:scope-dc-energy}
\end{equation}
와 fixed-cell nuclear gradient
\begin{equation}
  g_{Ax}
  = \frac{\dd E_{\mathrm{DC-DFA}}}{\dd R_{Ax}}
  \label{eq:scope-gradient}
\end{equation}
이다. $A$는 원자, $x\in\{x,y,z\}$는 Cartesian component다. Public API가 반환하는
값은 gradient이며 force는
\begin{equation}
  F_{Ax}=-g_{Ax}
  \label{eq:scope-force-sign}
\end{equation}
이다.

식~\eqref{eq:scope-dc-energy}는 ``HF energy에 DFA correction을 더한 식''이라는 뜻이
아니다. Target DFA의 kinetic, electron--nuclear, Hartree, XC, ion--ion energy를 모두
HF density에서 새로 평가한다. HF 계산의 energy는 source orbital을 만드는 reference와
diagnostic으로만 남는다.

\subsection{왜 ordinary KS gradient와 다른가}

Self-consistent KS-DFA에서는 orbital variation $\kappa$에 대해 같은 DFA energy가
stationary하다. 반면 DC-DFA에서는 orbital을 정하는 functional과 energy를 평가하는
functional이 다르므로
\begin{equation}
  \left.
  \frac{\partial E_{\mathrm{DC-DFA}}}{\partial\kappa_{ai}^{\vk}}
  \right|_{D_{\mathrm{HF}}}
  \ne 0
  \label{eq:scope-nonstationary}
\end{equation}
이다. 따라서 핵이 움직일 때 HF orbital이 변하는 효과를 무시할 수 없다. 모든 핵
변위마다 CPHF를 직접 푸는 대신 HF orbital Hessian의 adjoint equation, 즉 Z-vector
equation을 한 번 풀어 response를 모은다.

\subsection{Active 범위}

이 판에서 하나의 완결된 이론 대상으로 다루는 범위는 다음과 같다.

\begin{itemize}
  \item restricted closed-shell RHF/KRHF source density,
  \item three-dimensional periodic boundary condition,
  \item GTH pseudopotential과 atom-centered Gaussian AO,
  \item Gamma-point 및 일반 sampled k-point mesh,
  \item PBE, SCAN, r2SCAN target DFA,
  \item FFTDF, GDF, RSDF two-electron representation,
  \item CPU와 GPU implementation,
  \item fixed lattice에서의 energy 및 nuclear analytic gradient.
\end{itemize}

Stress/cell derivative, target-only multigrid, larger-angular-momentum completion
matrix와 open-shell/UHF response는 active 범위에서 제외되어 별도 archive에 보존된다.
따라서 이 문서에서 ``analytic gradient''라는 말은 항상 cell을 고정한
$\partial E/\partial R_{Ax}$를 뜻한다.

\subsection{필요한 사전 지식과 장 구성}

독자는 AO 기반 HF와 KS-DFT, density matrix, Gaussian integral, reciprocal lattice의
기초를 알고 있다고 가정한다. 제~\ref{sec:notation}장은 periodic notation과 단위를,
제~\ref{sec:fixed-density-energy}장은 fixed-HF-density target energy를 정의한다.
제~\ref{sec:orbital-response}장과 제~\ref{sec:dc-lagrangian}장은 nonvariational
response와 전체 gradient를 유도한다. 이후 GTH, FFTDF, GDF, MGGA, RSDF 순으로 각
explicit derivative를 설명하고 마지막에 검증 방법을 정리한다.

수식에서 ``정확하다''는 말은 두 층위를 구분해야 한다. 첫째, 선택한 finite basis,
mesh, auxiliary basis와 k-point representation에서 analytic derivative가 같은
representation의 full-energy finite difference와 일치해야 한다. 둘째, 그 finite
representation 자체가 cutoff, basis와 k-mesh에 대해 수렴해야 한다. 첫 번째는 구현
정확도이고 두 번째는 물리적 수렴도다.

\section{Periodic notation, matrix convention과 단위}
\label{sec:notation}

\subsection{Lattice와 reciprocal lattice}

Primitive lattice matrix를 column convention으로
\begin{equation}
  \mathbf A=(\mathbf a_1\;\mathbf a_2\;\mathbf a_3)
\end{equation}
라 하고 cell volume을
\begin{equation}
  \Omega=\det\mathbf A>0
  \label{eq:cell-volume}
\end{equation}
로 둔다. Reciprocal lattice matrix는
\begin{equation}
  \mathbf B=2\pi\mathbf A^{-\Herm},
  \qquad
  \mathbf a_i\cdot\mathbf b_j=2\pi\delta_{ij}
  \label{eq:reciprocal-lattice}
\end{equation}
이다. Lattice translation은 $\vT=\mathbf A\mathbf n$,
$\mathbf n\in\mathbb Z^3$이고 reciprocal vector는 $\vG=\mathbf B\mathbf m$다.

Cell 안의 적분은 $\int_\Omega\dd\mathbf r$로 표기한다. Fourier pair는 implementation의
FFT normalization에 따라 상수 위치가 달라질 수 있으므로, 이 문서에서는
\begin{equation}
  \rho(\vG)=\int_\Omega e^{-i\vG\cdot\mathbf r}\rho(\mathbf r)\dd\mathbf r,
  \qquad
  \rho(\mathbf r)=\frac{1}{\Omega}\sum_{\vG}
  e^{i\vG\cdot\mathbf r}\rho(\vG)
  \label{eq:fourier-pair}
\end{equation}
를 사용한다.

\subsection{Bloch AO와 k-point weight}

원자 중심 Gaussian $\chi_\mu$로부터 Bloch AO를
\begin{equation}
  \phi_{\mu\vk}(\mathbf r)
  =\frac{1}{\sqrt{N_{\mathrm c}}}
   \sum_{\vT}e^{i\vk\cdot\vT}
   \chi_\mu(\mathbf r-\vT)
  \label{eq:bloch-ao}
\end{equation}
로 정의한다. $N_{\mathrm c}$는 Born--von Karman supercell의 cell 수이며 최종 intensive
quantity에서는 소거된다. Bloch orbital은
\begin{equation}
  \psi_{p\vk}(\mathbf r)
  =\sum_\mu C_{\mu p}^{\vk}\phi_{\mu\vk}(\mathbf r)
  \label{eq:bloch-mo}
\end{equation}
이다.

일반 weight $w_\vk$는 $\sum_\vk w_\vk=1$로 normalization한다. 이 package의 regular
mesh는 equal weight를 사용하므로 $w_\vk=1/N_k$다. 아래 식에서 $1/N_k$가 보이면
equal-weight convention을 뜻한다. Gamma-only도 별도 scalar theory가 아니라
$N_k=1$, $\vk=\mathbf0$인 같은 tensor convention으로 다룬다.

AO overlap과 one-electron matrix는
\begin{align}
  S_{\mu\nu}^{\vk}
  &=\langle\phi_{\mu\vk}|\phi_{\nu\vk}\rangle,\\
  h_{\mu\nu}^{\vk}
  &=\langle\phi_{\mu\vk}|\hat h|\phi_{\nu\vk}\rangle
  \label{eq:ao-matrices}
\end{align}
이고 MO coefficient는
\begin{equation}
  C^{\vk\dagger}S^{\vk}C^{\vk}=I
  \label{eq:mo-orthonormality}
\end{equation}
를 만족한다.

\subsection{Restricted density matrix}

Closed-shell occupation $f_{p\vk}\in\{0,2\}$를 사용하면 AO density는
\begin{equation}
  D_{\mu\nu}^{\vk}
  =\sum_p f_{p\vk}C_{\mu p}^{\vk}C_{\nu p}^{\vk *}
  =2\sum_i C_{\mu i}^{\vk}C_{\nu i}^{\vk *}
  \label{eq:restricted-density}
\end{equation}
이다. 이 정의는 occupation factor $2$를 이미 포함한다. 따라서 energy와 response
contraction에서 factor $2$를 다시 넣는지 여부는 사용한 occupied-only tensor가
$D$를 대체하는지에 따라 결정해야 한다.

Matrix trace convention은
\begin{equation}
  \Tr[D^{\vk}X^{\vk}]
  =\sum_{\mu\nu}D_{\nu\mu}^{\vk}X_{\mu\nu}^{\vk}
  \label{eq:trace-convention}
\end{equation}
이다. Complex k-point에서 $D^{\vk}$와 Hermitian operator는 Hermitian이지만 개별
derivative block은 반드시 Hermitian일 필요가 없다. 최종 scalar에는 실수부를 취하되,
중간 단계의 imaginary part를 임의로 버리지 않는다.

\subsection{Index와 derivative notation}

\begin{center}
\begin{tabular}{cl}
\toprule
Index & 의미 \\
\midrule
$\mu,\nu,\lambda,\sigma$ & primary AO \\
$P,Q$ & auxiliary AO \\
$i,j$ & occupied MO \\
$a,b$ & virtual MO \\
$p,q$ & arbitrary MO \\
$A,B$ & atom \\
$x,y$ & Cartesian component \\
$\vk,\vk'$ & sampled k-point \\
\bottomrule
\end{tabular}
\end{center}

핵 좌표 $R_{Ax}$에 대한 total derivative는 superscript $Ax$ 또는
$\partial_{Ax}$로 쓴다. 예를 들어 $h^{Ax}=\partial h/\partial R_{Ax}$다. AO basis가
핵과 함께 움직이므로 matrix derivative와 orbital coefficient response를 구분해야 한다.
Cell은 고정되어 있어 $\mathbf A^{Ax}=0$이고 reciprocal vectors와 discrete mesh도
고정된다.

\subsection{Canonical unit}

\begin{center}
\begin{tabular}{lll}
\toprule
Quantity & Canonical unit & 비고 \\
\midrule
Energy & Hartree (Ha) & public energy \\
Position/lattice & Bohr & PySCF internal \\
Nuclear gradient & Ha/Bohr & $\dd E/\dd R$ \\
Force & Ha/Bohr & $-\dd E/\dd R$ \\
k-point & Bohr$^{-1}$ & Cartesian reciprocal coordinate \\
Plane-wave cutoff & Ha & $G_{\max}^2/2$ \\
\bottomrule
\end{tabular}
\end{center}

입력 XYZ가 Angstrom이어도 cell을 build한 뒤 내부 좌표는 Bohr다. Gradient array를
eV/Angstrom으로 내보낼 때는 energy와 length conversion을 동시에 적용하며, 부호를
바꾸는 것은 unit conversion이 아니라 force 변환이다.

\section{HF density에서의 target DFA energy}
\label{sec:fixed-density-energy}

\subsection{Source HF problem}

Source orbital은 periodic RHF/KRHF equation
\begin{equation}
  F_{\mathrm{HF}}^{\vk}C^{\vk}
  =S^{\vk}C^{\vk}\varepsilon^{\vk}
  \label{eq:hf-roothaan}
\end{equation}
을 self-consistent하게 풀어 얻는다. Restricted HF Fock은
\begin{equation}
  F_{\mathrm{HF}}^{\vk}
  =h^{\vk}+J^{\vk}[D]-\frac12K^{\vk}[D]
  \label{eq:hf-fock}
\end{equation}
이다. $D$는 occupation factor $2$를 포함한 식~\eqref{eq:restricted-density}의 density다.
HF energy per primitive cell은
\begin{equation}
  E_{\mathrm{HF}}
  =\frac{1}{N_k}\sum_{\vk}\Tr[D^{\vk}h^{\vk}]
   +\frac{1}{2N_k}\sum_{\vk}\Tr\!\left[
      D^{\vk}\left(J^{\vk}-\frac12K^{\vk}\right)\right]
   +E_{\mathrm{NN}}
  \label{eq:hf-energy}
\end{equation}
이다. Source HF 계산에서 FFTDF, GDF 또는 RSDF는 $J$와 $K$의 numerical
representation을 정한다.

\subsection{Frozen-HF-density target energy}

Source HF가 수렴하면 density를 고정한 채 target RKS/KRKS object를 만들고 target
functional의 energy를 평가한다. Semilocal target DFA에서
\begin{equation}
  E_{\mathrm{DC-DFA}}[D]
  =E_1[D]+E_{\mathrm H}[D]+E_{\mathrm{xc}}^{\mathrm{DFA}}[D]
   +E_{\mathrm{NN}}
  \label{eq:dc-energy-components}
\end{equation}
이며
\begin{align}
  E_1[D]
  &=\frac{1}{N_k}\sum_{\vk}\Tr[D^{\vk}h^{\vk}],
  \label{eq:dc-one-electron-energy}\\
  E_{\mathrm H}[D]
  &=\frac{1}{2N_k}\sum_{\vk}\Tr[D^{\vk}J^{\vk}[D]].
  \label{eq:dc-hartree-energy}
\end{align}
Target PBE, SCAN, r2SCAN은 pure semilocal functional이므로 target energy에 exact
exchange $K$는 없다. 그러나 $D$가 HF solution이므로 orbital response를 정하는
Hessian에는 source HF exchange가 반드시 남는다.

Target Fock, 즉 $E_{\mathrm{DC-DFA}}$의 density derivative는
\begin{equation}
  F_{\mathrm t}^{\vk}
  =h^{\vk}+J^{\vk}[D]+V_{\mathrm{xc}}^{\vk}[D]
  \label{eq:target-fock}
\end{equation}
이다. 여기서 subscript $\mathrm t$는 target을 뜻한다. Energy evaluation에서
$D_{\mathrm{HF}}$를 넘긴 뒤 target KS orbital optimization을 수행하지 않는다.

\subsection{Density와 semilocal variables}

Real-space density는
\begin{equation}
  \rho(\mathbf r)
  =\frac{1}{N_k}\sum_{\vk}\sum_{\mu\nu}
    D_{\mu\nu}^{\vk}
    \phi_{\nu\vk}^{*}(\mathbf r)
    \phi_{\mu\vk}(\mathbf r)
  \label{eq:real-space-density}
\end{equation}
이다. Spin-unpolarized GGA는
\begin{equation}
  E_{\mathrm{xc}}^{\mathrm{GGA}}
  =\int_\Omega
    f\!\left(\rho,\sigma\right)\dd\mathbf r,
  \qquad
  \sigma=|\nabla\rho|^2
  \label{eq:gga-energy}
\end{equation}
로 쓸 수 있다.

SCAN과 r2SCAN은 positive kinetic-energy density
\begin{equation}
  \tau(\mathbf r)
  =\frac{1}{2N_k}\sum_{\vk}\sum_{\mu\nu}
   D_{\mu\nu}^{\vk}
   \nabla\phi_{\nu\vk}^{*}(\mathbf r)\cdot
   \nabla\phi_{\mu\vk}(\mathbf r)
  \label{eq:tau-definition}
\end{equation}
에도 의존하여
\begin{equation}
  E_{\mathrm{xc}}^{\mathrm{MGGA}}
  =\int_\Omega f(\rho,\sigma,\tau)\dd\mathbf r
  \label{eq:mgga-energy}
\end{equation}
가 된다. PySCF의 SCAN/r2SCAN route는 density Laplacian channel을 사용하지 않는다.
식~\eqref{eq:tau-definition}의 $1/2$는 convention의 일부이므로 $v_\tau=\partial
f/\partial\tau$ contraction에서 정확히 한 번 반영해야 한다.

\subsection{Nonvariational orbital gradient}

MO coefficient를 occupied--virtual anti-Hermitian rotation으로 변화시키자.
\begin{equation}
  C^{\vk}\longrightarrow C^{\vk}e^{\kappa^{\vk}},
  \qquad \kappa^{\vk\dagger}=-\kappa^{\vk}.
  \label{eq:orbital-rotation}
\end{equation}
Closed-shell energy의 first variation은 convention에 따라 factor $2$를 포함해
\begin{equation}
  \delta E_{\mathrm{DC-DFA}}
  =\frac{2}{N_k}\operatorname{Re}
   \sum_{\vk ai}F_{\mathrm t,ai}^{\vk}\,
   \delta\kappa_{ai}^{\vk}
  \label{eq:target-orbital-gradient}
\end{equation}
로 쓸 수 있다. Source HF canonical orbital은
$F_{\mathrm{HF},ai}^{\vk}=0$을 만족하지만 일반적으로
\begin{equation}
  F_{\mathrm t,ai}^{\vk}\ne0.
  \label{eq:target-vo-nonzero}
\end{equation}
이 nonzero target VO block이 Z-vector right-hand side다.

\subsection{Energy object가 보존해야 할 metadata}

같은 이름의 DFA라도 basis, pseudopotential, cell, k-point, DF backend와 numerical
grid가 다르면 서로 다른 finite energy다. 재현 가능한 energy object는 최소한 target
functional, source reference, $N_k$, lattice와 coordinate unit, primary/auxiliary basis,
GTH pseudopotential, DF backend, FFT mesh 또는 RSDF source/metric parameter를 함께
기록해야 한다. Gradient는 이 energy object와 동일한 설정을 미분해야 한다.

\section{Periodic orbital response와 Z-vector}
\label{sec:orbital-response}

\subsection{Orthonormality response}

핵 변위 $R_{Ax}$에서 AO overlap이 변하므로 MO coefficient derivative는 단순한
unitary rotation만이 아니다. 식~\eqref{eq:mo-orthonormality}을 미분하면
\begin{equation}
  C^{\vk\dagger}S^{\vk}C^{\vk,Ax}
  +C^{\vk,Ax\dagger}S^{\vk}C^{\vk}
  +C^{\vk\dagger}S^{\vk,Ax}C^{\vk}=0
  \label{eq:orthonormality-response}
\end{equation}
이다. $U^{\vk,Ax}=C^{\vk\dagger}S^{\vk}C^{\vk,Ax}$라 두면
\begin{equation}
  U^{\vk,Ax}+U^{\vk,Ax\dagger}
  =-S_{\mathrm{MO}}^{\vk,Ax},
  \qquad
  S_{\mathrm{MO}}^{\vk,Ax}=C^{\vk\dagger}S^{\vk,Ax}C^{\vk}
  \label{eq:u-overlap-constraint}
\end{equation}
이다. Occupied--occupied와 virtual--virtual gauge는 선택할 수 있지만 occupied--virtual
block은 HF response equation으로 결정된다.

\subsection{CPHF equation}

HF stationarity condition $F_{ai}^{\vk}=0$을 핵 좌표로 미분하면
\begin{equation}
  \sum_{\vk'bj}A_{ai\vk,bj\vk'}U_{bj}^{\vk',Ax}
  =-B_{ai}^{\vk,Ax}
  \label{eq:cphf-equation}
\end{equation}
을 얻는다. $A$는 periodic restricted HF orbital Hessian이고 $B^{Ax}$는 nuclear
perturbation source다. Canonical-orbital notation에서 Hessian action은 개략적으로
\begin{equation}
  (AU)_{ai}^{\vk}
  =(\varepsilon_a^{\vk}-\varepsilon_i^{\vk})U_{ai}^{\vk}
  +\left[C_v^{\vk\dagger}\,\delta G_{\mathrm{HF}}[D^{(1)}(U)]^{\vk}
  C_o^{\vk}\right]_{ai}
  \label{eq:hf-hessian-action}
\end{equation}
이다. $\delta G_{\mathrm{HF}}=\delta J-\tfrac12\delta K$이며, 따라서 선택한
FFTDF/GDF/RSDF의 transition-density J/K가 AX operation의 핵심이다.

Nuclear source는 one-electron, two-electron과 overlap 항을 함께 포함한다.
\begin{equation}
  B_{ai}^{\vk,Ax}
  =\left[C_v^{\vk\dagger}
    \left(h^{\vk,Ax}+G_{\mathrm{HF}}^{\vk,Ax}
          -\varepsilon_i^{\vk}S^{\vk,Ax}\right)
    C_o^{\vk}\right]_{ai}
  +B_{ai,\mathrm{basis}}^{\vk,Ax}.
  \label{eq:cphf-source}
\end{equation}
$G_{\mathrm{HF}}^{Ax}$는 density를 고정한 integral derivative이고 마지막 항은 chosen
AO-row/column convention에 따라 식의 다른 항에 흡수될 수 있다. Production code와
문서가 같은 convention을 사용해야 한다.

\subsection{Adjoint Z-vector}

모든 $3N_{\mathrm{atom}}$ 변위에 대해 식~\eqref{eq:cphf-equation}을 풀면 비용이 크다.
Target orbital gradient를 vector $g$로 모아 adjoint equation
\begin{equation}
  A^{\Herm}z=-g
  \label{eq:zvector-equation}
\end{equation}
을 한 번 푼다. 그러면 chain-rule contribution은
\begin{equation}
  g_{Ax}^{\mathrm{resp}}
  =z^{\Herm}B^{Ax}
  \label{eq:zvector-contraction}
\end{equation}
으로 계산된다. 식~\eqref{eq:zvector-equation}과
\eqref{eq:zvector-contraction}의 부호는 CPHF equation을 $AU=-B$로 정의한 것과
연동된다. $B$ 또는 $g$ 정의를 바꾸면 두 식의 부호를 함께 바꿔야 한다.

\subsection{Complex k-point의 real-linear 문제}

Complex orbital rotation $U=X+iY$가 만드는 Hermitian first-order density는
\begin{equation}
  D^{(1),\vk}(U)
  =2\left(C_v^{\vk}U^{\vk}C_o^{\vk\dagger}
          +C_o^{\vk}U^{\vk\dagger}C_v^{\vk\dagger}\right)
  \label{eq:first-order-density}
\end{equation}
이다. $U^*$가 동시에 나타나므로 map $U\mapsto D^{(1)}$는 complex-linear가 아니라
real-linear다. 따라서 solver coordinate를
\begin{equation}
  u_R=\begin{pmatrix}\operatorname{Re}U\\\operatorname{Im}U\end{pmatrix}
  \label{eq:realified-coordinate}
\end{equation}
로 두고
\begin{equation}
  A_Ru_R=-B_R,
  \qquad A_R^Tz_R=-g_R
  \label{eq:realified-zvector}
\end{equation}
를 푼다. Gamma에서 matrix가 실수라는 사실은 이 일반식을 축약하는 특수 경우다.

\subsection{k-point normalization과 momentum structure}

Nuclear perturbation은 primitive-cell periodicity를 유지하므로 같은 $\vk$ 안에서 VO
response를 만든다. 그러나 $J/K$ Hessian action은 모든 k-point density response를
couple한다. Equal-weight convention에서 AX와 최종 scalar contraction 모두 정확히
$1/N_k$를 한 번 포함해야 한다. Transition density가 occupation factor를 이미 포함하는지
확인하지 않고 별도의 factor $2$를 곱하면 Gamma에서는 우연히 맞아도 sampled k-point에서
일관성이 깨질 수 있다.

\subsection{Gauge, 작은 gap과 residual}

Occupied 또는 virtual subspace 안의 unitary rotation은 density를 바꾸지 않는다.
Degenerate orbital에서는 개별 $U_{ai}$보다 subspace-invariant density response가
물리량이다. Solver는 VO 공간만 사용하고, denominator가 작은 경우 regularization을
숨겨 적용하지 않는다. Report에는 최소한
\begin{equation}
  r_{\max}=\|A_R^Tz_R+g_R\|_\infty
  \label{eq:zvector-residual}
\end{equation}
와 response dimension, condition estimate를 기록한다. Gradient--FD agreement가 좋아도
$r_{\max}$가 크면 특정 perturbation에서 cancellation에 의존할 수 있으므로 완료 근거로
충분하지 않다.

\section{Periodic density fitting과 RSDF response}
\label{sec:periodic-df}

\subsection{왜 density fitting의 미분이 필요한가}

Periodic HF와 target DFA의 Coulomb 및 exact-exchange energy는 four-center electron
repulsion integral을 직접 다루면 계산량과 memory가 빠르게 증가한다. Density
fitting은 AO pair density를 auxiliary Gaussian basis로 전개하여 이 문제를
three-center quantity와 two-center metric으로 바꾼다. 그러나 핵이 움직이면 primary
AO와 auxiliary basis가 모두 함께 움직인다. 따라서 energy value에 density fitting을
사용했다면 gradient에도 auxiliary basis와 metric response가 포함되어야 한다.

Bloch AO를 $\chi_{\mu\vk}(\mathbf r)$, momentum transfer를
\begin{equation}
 \mathbf q=\mathbf k_j-\mathbf k_i
\end{equation}
로 정의한다. Auxiliary Bloch function $P\mathbf q$에 대한 three-center source와
Coulomb metric은
\begin{align}
 B_{\mu\nu,P}^{ij}
 &=\left(\chi_{\mu\mathbf k_i}\chi_{\nu\mathbf k_j}\middle|P\mathbf q\right),\\
 M_{PQ}(\mathbf q)
 &=\left(P\mathbf q\middle|Q\mathbf q\right)
\end{align}
이다. 여기서 ket AO에 complex conjugate가 포함되는지 여부는 program의 integral
orientation convention에 따라 달라질 수 있다. 중요한 점은 central value와
derivative가 같은 convention을 사용해야 한다는 것이다.

Positive-definite metric의 Cholesky decomposition을
\begin{equation}
 \mathbf M(\mathbf q)=\mathbf L(\mathbf q)\mathbf L(\mathbf q)^\dagger
\end{equation}
로 두면 orthonormalized three-center factor는
\begin{equation}
 \mathbf C^{ij}=\mathbf L(\mathbf q)^{-1}\mathbf B^{ij}
 \label{eq:cderi-definition}
\end{equation}
이다. Four-center integral은 대략
\begin{equation}
 (\mu i,\nu j|\lambda l,\sigma m)
 \simeq \sum_P C_{P,\mu\nu}^{ij}
 C_{P,\lambda\sigma}^{lm*}
\end{equation}
처럼 factorized된다. $\mathbf C$가 PySCF의 CDERI에 해당한다.

\subsection{Nuclear derivative와 auxiliary response}

Atom $A$의 Cartesian coordinate $R_{Ax}$를 미분하면 식
\eqref{eq:cderi-definition}에서
\begin{equation}
 \mathbf C^{Ax}
 =\mathbf L^{-1}\mathbf B^{Ax}
 -\mathbf L^{-1}\mathbf L^{Ax}\mathbf C
 \label{eq:cderi-response}
\end{equation}
를 얻는다. 첫 항은 three-center source response이고 둘째 항은 metric response이다.
Metric을 고정하면 auxiliary Gaussian center가 움직이는 효과와 fitting coefficient의
변화를 잃는다. 이 누락은 분자가 작을 때 우연히 작을 수 있지만 auxiliary center가
많아질수록 일반적으로 커진다.

Cholesky derivative는
\begin{equation}
 \mathbf M^{Ax}=\mathbf L^{Ax}\mathbf L^\dagger
 +\mathbf L\mathbf L^{Ax\dagger}
 \label{eq:cholesky-response}
\end{equation}
를 만족한다. 양쪽에 $\mathbf L^{-1}$와 $\mathbf L^{-\dagger}$를 곱하여
\begin{equation}
 \mathbf W=\mathbf L^{-1}\mathbf M^{Ax}\mathbf L^{-\dagger}
\end{equation}
를 정의한다. Lower-Cholesky gauge에서는
\begin{equation}
 \mathbf L^{Ax}=\mathbf L\,\Phi(\mathbf W)
\end{equation}
로 쓸 수 있다. $\Phi$는 strict lower triangle을 보존하고 Hermitian diagonal의
절반을 취한다. 이 gauge를 central value와 derivative에서 일관되게 사용해야 한다.

\subsection{Coulomb kernel의 range separation}

RSDF는 Coulomb kernel을
\begin{equation}
 \frac{1}{r}
 =\frac{\operatorname{erfc}(\omega r)}{r}
 +\frac{\operatorname{erf}(\omega r)}{r}
 \label{eq:range-separation}
\end{equation}
로 나눈다. 첫 항은 real space에서 빠르게 감쇠하므로 lattice sum으로 계산하고, 둘째
항은 reciprocal space의 smooth Fourier sum으로 계산한다. 이 분해는 exact identity이지만
finite image list와 finite reciprocal mesh를 사용하면 $\omega$, screening precision,
mesh가 실제 numerical representation의 일부가 된다.

PySCF public RSDF에서는 three-center source와 two-center metric의 최적 parameter가
독립적으로 결정될 수 있다. 이를 명시하면
\begin{align}
 \mathbf B^{ij}
 &=\mathbf B_{\mathrm{SR}}^{ij}(\omega_B,\epsilon_R)
 +\mathbf B_{\mathrm{LR}}^{ij}(\omega_B,\mathcal G_B)
 +\delta_{\mathbf q,0}\mathbf B_{vbar}^{ij},
 \label{eq:rsdf-source}\\
 \mathbf M(\mathbf q)
 &=\mathbf M_{\mathrm{SR}}(\mathbf q;\omega_M,\epsilon_M)
 +\mathbf M_{\mathrm{LR}}(\mathbf q;\omega_M,\mathcal G_M)
 +\delta_{\mathbf q,0}\mathbf M_{G=0}.
 \label{eq:rsdf-metric}
\end{align}
일반적으로 $\omega_B\neq\omega_M$이고 $\mathcal G_B\neq\mathcal G_M$이다. Auxiliary
basis의 exponent 분포는 primary AO pair의 분포와 다르므로 metric에 더 높은 cutoff가
필요할 수 있다.

핵 변위에서 fixed cell을 가정하면 lattice vector와 reciprocal vector는 변하지 않는다.
따라서 $\omega$와 mesh를 numerical input으로 고정한 derivative는
\begin{align}
 \mathbf B^{ij,Ax}
 &=\mathbf B_{\mathrm{SR}}^{ij,Ax}
 +\mathbf B_{\mathrm{LR}}^{ij,Ax}
 +\delta_{\mathbf q,0}\mathbf B_{vbar}^{ij,Ax},\\
 \mathbf M^{Ax}
 &=\mathbf M_{\mathrm{SR}}^{Ax}
 +\mathbf M_{\mathrm{LR}}^{Ax}
 +\delta_{\mathbf q,0}\mathbf M_{G=0}^{Ax}
\end{align}
이다. Stress에서는 lattice와 reciprocal mesh 자체가 변하므로 이 장의 fixed-cell 식만으로
충분하지 않다.

\subsection{Primary AO와 auxiliary-center ownership}

Three-center integral에는 세 개의 atom-centered object가 있다. Symbolically
\begin{equation}
 B_{\mu\nu,P}^{Ax}
 =\delta_{A,A_\mu}\partial_{A_\mu x}B_{\mu\nu,P}
 +\delta_{A,A_\nu}\partial_{A_\nu x}B_{\mu\nu,P}
 +\delta_{A,A_P}\partial_{A_Px}B_{\mu\nu,P}.
\end{equation}
Integral library의 \texttt{ip1}, \texttt{ip2}는 흔히 electronic coordinate에 대한 derivative를
반환한다. Nuclear center를 이동시키는 derivative에는 translational identity 때문에
반대 부호가 붙는다. AO row, AO column, auxiliary index 각각의 atom slice를 적용한 뒤
세 contribution을 합해야 한다.

Two-center metric에는 두 auxiliary center가 있으므로
\begin{equation}
 M_{PQ}^{Ax}
 =\delta_{A,A_P}\partial_{A_Px}M_{PQ}
 +\delta_{A,A_Q}\partial_{A_Qx}M_{PQ}.
\end{equation}
Short-range metric derivative는 central value와 동일한 screened lattice image list를
사용해야 한다. Image cutoff가 다르면 둘 다 연속적인 Coulomb metric을 근사하더라도
finite calculation의 exact derivative가 아니다.

\subsection{\texorpdfstring{$G=0$}{G=0} mode와 charge correction}

Periodic Coulomb kernel의 $G=0$ mode는 charge neutrality와 Ewald convention에 의해
특별히 처리된다. Source의 \texttt{vbar}와 metric의 $G=0$ correction은 momentum transfer가
정확히 zero인 block에서만 나타난다.
\begin{equation}
 \delta_{\mathbf q,0}=\begin{cases}
 1,& \mathbf k_j-\mathbf k_i=\mathbf 0,\\
 0,& \text{otherwise}.
 \end{cases}
\end{equation}
Gamma-only calculation뿐 아니라 sampled k-point calculation의 diagonal pair도
$\mathbf q=0$ block이다. 따라서 ``Gamma인가?''와 ``momentum transfer가 zero인가?''를
동일한 조건으로 취급하면 안 된다.

\subsection{k-point pair와 Hermitian reconstruction}

$N_k$개의 k-point가 있으면 J/K build에는 여러 $(i,j)$ pair가 필요하다. Momentum
conservation 때문에 metric은 $\mathbf q=\mathbf k_j-\mathbf k_i$에 의해 분류된다.
Canonical pair만 저장할 때 reverse pair는
\begin{equation}
 C^{ji}_{P,\nu\mu}=C^{ij*}_{P,\mu\nu}
\end{equation}
와 같은 Hermitian relation으로 복원한다. Nuclear derivative도 같은 conjugation과 AO
index transpose를 따라야 한다. Complex k-point에서 실수 gradient를 얻으려면 모든
pair contraction과 equal-weight $1/N_k$ normalization이 일관되어야 한다.

Closed-shell Coulomb auxiliary density를
\begin{equation}
 j_P=\sum_{\vk\mu\nu}C^\vk_{P,\mu\nu}D^\vk_{\nu\mu}
\end{equation}
라 하면 fixed-density J derivative는
\begin{equation}
 (J^\vk)^{Ax}_{\mu\nu}
 =\frac{1}{N_k}\sum_P
 \left[j_P^{Ax}C^\vk_{P,\mu\nu}
 +j_PC^{\vk,Ax}_{P,\mu\nu}\right].
\end{equation}
Exchange에는 두 CDERI factor가 있으므로 두 factor 각각에 product rule을 적용한다.

\subsection{DC-DFT gradient 안에서의 역할}

RSDF는 two-electron value와 response의 numerical representation을 정한다. DC-DFT의
nonvariational 구조 자체는 backend에 의존하지 않는다. 전체 gradient는
\begin{equation}
 E_{\mathrm{DC}}^{Ax}
 =E_{\mathrm{target,explicit}}^{Ax}
 +\mathbf z^T\mathbf B_{\mathrm{HF}}^{Ax}
 \label{eq:dc-rsdf-total}
\end{equation}
로 쓸 수 있다. 첫 항에는 target one-electron, Hartree, XC, overlap/Pulay 및 nuclear
contribution이 들어간다. 둘째 항은 HF density가 핵 변위에 따라 변하는 orbital
response를 adjoint Z-vector로 포함한다. RSDF의 $C^{Ax}$는 target explicit J와 source
HF response operator 양쪽에 사용된다.

\subsection{검증 원리}

Analytic derivative의 correctness와 physical convergence는 다른 문제이다. 먼저 같은
finite RSDF representation에서 central energy finite difference와 비교한다.
\begin{equation}
 g_{Ax}^{\mathrm{FD}}(h)
 =\frac{E(\mathbf R+h\mathbf e_{Ax})
       -E(\mathbf R-h\mathbf e_{Ax})}{2h}.
\end{equation}
그 다음 $h$, source/metric mesh, primary/auxiliary basis, k-mesh를 각각 수렴시킨다.

또한 독립적으로 재구성한 CDERI가 SCF storage와 일치해야 한다.
\begin{equation}
 \Delta_C=\max_{ijP\mu\nu}
 \left|C_{P,\mu\nu}^{ij,\mathrm{reconstructed}}
 -C_{P,\mu\nu}^{ij,\mathrm{storage}}\right|.
\end{equation}
$\Delta_C$가 작지 않으면 derivative 식이 옳더라도 다른 omega, mesh, pair orientation,
또는 factor gauge를 미분한 것이므로 production gradient를 거부해야 한다.

\section{DC-DFT Lagrangian과 전체 nuclear gradient}
\label{sec:dc-lagrangian}

\subsection{Constraint를 포함한 Lagrangian}

HF orbital coefficient를 $C$라 하고 HF stationarity residual의 VO block을
\begin{equation}
  f_{ai}^{\vk}(C,\vR)
  =\left[C_v^{\vk\dagger}F_{\mathrm{HF}}^{\vk}C_o^{\vk}\right]_{ai}
  \label{eq:hf-stationarity-residual}
\end{equation}
로 정의한다. 수렴한 HF solution에서는 $f=0$이다. Target energy 자체는 이 constraint에
대해 stationary하지 않으므로 adjoint multiplier $z$를 도입한 reduced Lagrangian을
\begin{equation}
  \mathcal L(C,z;\vR)
  =E_{\mathrm{DC-DFA}}(C;\vR)
   +2\operatorname{Re}\sum_{\vk ai}w_\vk
    z_{ai}^{\vk *}f_{ai}^{\vk}(C;\vR)
  \label{eq:dc-lagrangian}
\end{equation}
로 둔다. MO orthonormality multiplier는 아래의 energy-weighted density와 overlap/Pulay
항에 흡수하여 표기한다.

Orbital rotation에 대한 stationarity
\begin{equation}
  \frac{\partial\mathcal L}{\partial\kappa}=0
  \label{eq:lagrangian-orbital-stationarity}
\end{equation}
를 적용하면 식~\eqref{eq:zvector-equation}의 Z-vector equation을 얻는다. 따라서 total
derivative에서 명시적인 $\dd C/\dd R_{Ax}$ 항은 사라지고
\begin{equation}
  \frac{\dd E_{\mathrm{DC-DFA}}}{\dd R_{Ax}}
  =\left.\frac{\partial\mathcal L}{\partial R_{Ax}}\right|_{C,z}
  =g_{Ax}^{\mathrm{explicit}}+g_{Ax}^{\mathrm{response}}
  \label{eq:dc-gradient-master}
\end{equation}
만 남는다.

\subsection{Explicit target contribution}

Density matrix를 고정한 target energy의 explicit derivative를 다음처럼 분해한다.
\begin{equation}
  g_{Ax}^{\mathrm{explicit}}
  =g_{Ax}^{1e}+g_{Ax}^{\mathrm H}+g_{Ax}^{\mathrm{xc}}
   +g_{Ax}^{\mathrm{Pulay,t}}+g_{Ax}^{\mathrm{NN}}.
  \label{eq:explicit-gradient-decomposition}
\end{equation}

One-electron term은
\begin{equation}
  g_{Ax}^{1e}
  =\frac{1}{N_k}\sum_\vk
    \Tr[D^{\vk}h^{\vk,Ax}]
  \label{eq:one-electron-gradient}
\end{equation}
이다. $h^{Ax}$에는 kinetic, local GTH, nonlocal GTH의 AO matrix derivative가 모두
포함된다.

Hartree term은 density를 고정한 Coulomb matrix derivative로
\begin{equation}
  g_{Ax}^{\mathrm H}
  =\frac{1}{2N_k}\sum_\vk
    \Tr[D^{\vk}J^{\vk,Ax}[D]]
  \label{eq:hartree-gradient}
\end{equation}
이다. $J^{Ax}[D]$는 FFTDF에서는 AO/grid와 Coulomb-kernel derivative, GDF/RSDF에서는
three-center source와 two-center metric derivative를 포함한다.

Semilocal XC explicit derivative는 density matrix를 고정하고 AO와 quadrature representation을
미분한 값이다.
\begin{equation}
  g_{Ax}^{\mathrm{xc}}
  =\left.
    \frac{\partial E_{\mathrm{xc}}[D;\vR]}{\partial R_{Ax}}
   \right|_D.
  \label{eq:xc-explicit-gradient}
\end{equation}
Atom-centered Gaussian AO를 uniform periodic grid에서 평가하므로 AO center derivative와
periodic image/Bloch phase convention이 central value와 같아야 한다. SCAN/r2SCAN에서는
$\tau$ channel이 추가된다.

\subsection{Target overlap/Pulay term}

Target Fock의 occupied--occupied block으로 target energy-weighted density를
\begin{equation}
  W_{\mathrm t}^{\vk}
  =2C_o^{\vk}
    \left(C_o^{\vk\dagger}F_{\mathrm t}^{\vk}C_o^{\vk}\right)
    C_o^{\vk\dagger}
  \label{eq:target-energy-weighted-density}
\end{equation}
로 정의한다. 그러면 target overlap contribution은
\begin{equation}
  g_{Ax}^{\mathrm{Pulay,t}}
  =-\frac{1}{N_k}\sum_\vk
    \Tr[W_{\mathrm t}^{\vk}S^{\vk,Ax}]
  \label{eq:target-overlap-pulay}
\end{equation}
이다. AO-row derivative formulation에서는 식~\eqref{eq:one-electron-gradient}--
\eqref{eq:target-overlap-pulay}의 일부가 atom-row contraction으로 재배치될 수 있지만,
합계는 동일해야 한다.

Ion--ion 항은 GTH ionic charge와 periodic Ewald convention으로
\begin{equation}
  g_{Ax}^{\mathrm{NN}}
  =\frac{\partial E_{\mathrm{NN}}}{\partial R_{Ax}}
  \label{eq:nuclear-repulsion-gradient}
\end{equation}
이다.

\subsection{HF orbital-response contribution}

식~\eqref{eq:zvector-contraction}을 index로 쓰면
\begin{equation}
  g_{Ax}^{\mathrm{response}}
  =2\operatorname{Re}\sum_{\vk ai}w_\vk
   z_{ai}^{\vk *}B_{ai}^{\vk,Ax}.
  \label{eq:response-gradient}
\end{equation}
여기서 $B^{Ax}$는 source HF equation의 nuclear derivative다. Target PBE와 MGGA가
바뀌어도 $B^{Ax}$와 HF Hessian의 물리적 정의는 바뀌지 않는다. 바뀌는 것은 target
VO gradient $g$와 explicit XC derivative이므로 같은 response infrastructure를 재사용할
수 있다.

\subsection{Backend-independent assembly}

최종 package decomposition은
\begin{align}
  g_{Ax}^{\mathrm{elec,explicit}}
  &=g_{Ax}^{1e}+g_{Ax}^{\mathrm H}+g_{Ax}^{\mathrm{xc}}
    +g_{Ax}^{\mathrm{Pulay,t}},\\
  g_{Ax}
  &=g_{Ax}^{\mathrm{elec,explicit}}
    +g_{Ax}^{\mathrm{NN}}+g_{Ax}^{\mathrm{response}}
  \label{eq:final-gradient-assembly}
\end{align}
로 기록한다. FFTDF, GDF, RSDF의 차이는 $J/K$ value와 derivative provider에 있고,
Z-vector equation 및 final assembly contract는 공통이다.

\begin{remark}
Full-energy finite difference는 $\vR\pm h\mathbf e_{Ax}$ 각각에서 source HF를 다시
수렴시키고 target energy를 새로 평가해야 한다. Central density를 고정한 component FD는
식~\eqref{eq:xc-explicit-gradient} 같은 개별 kernel 진단에는 유용하지만 전체 analytic
gradient의 reference가 아니다.
\end{remark}

\subsection{Translational invariance}

고정된 lattice 안에서 모든 원자를 같은 벡터만큼 이동하면 energy는 변하지 않아야 한다.
따라서 converged representation에서
\begin{equation}
  \sum_A\mathbf g_A=\mathbf0
  \label{eq:translation-sum-rule}
\end{equation}
이다. 큰 net gradient는 AO image derivative, auxiliary-center response, XC grid aliasing,
$G=0$ 처리 또는 k-point normalization 누락을 찾는 강력한 diagnostic이다.

\section{One-electron operator와 GTH pseudopotential derivative}
\label{sec:gth-one-electron}

\subsection{Periodic one-electron Hamiltonian}

GTH pseudopotential을 사용하는 one-electron Hamiltonian은
\begin{equation}
  \hat h
  =-\frac12\nabla^2
   +\sum_A\hat V_{\mathrm{loc}}^A
   +\sum_A\hat V_{\mathrm{nl}}^A
  \label{eq:gth-hamiltonian}
\end{equation}
이다. All-electron electron--nuclear $-Z_A/r$ 대신 smooth local potential과 finite-rank
nonlocal projector를 사용한다. Kinetic operator 자체는 핵 좌표에 명시적으로 의존하지
않지만 atom-centered AO basis가 움직이므로 kinetic AO matrix derivative는 0이 아니다.

\subsection{Local GTH potential}

대표적인 local form은
\begin{equation}
  V_{\mathrm{loc}}^A(r)
  =-\frac{Z_A^{\mathrm{ion}}}{r}
    \operatorname{erf}\!\left(\frac{r}{\sqrt2 r_{\mathrm{loc},A}}\right)
   +e^{-r^2/(2r_{\mathrm{loc},A}^2)}
    \sum_{n=1}^{4}C_{nA}
    \left(\frac{r}{r_{\mathrm{loc},A}}\right)^{2n-2}
  \label{eq:gth-local-real}
\end{equation}
이다. Periodic FFT representation에서는 원자 위치 dependence가 structure factor
$e^{-i\vG\cdot\mathbf R_A}$에 들어가므로
\begin{equation}
  \frac{\partial}{\partial R_{Ax}}
  \left[e^{-i\vG\cdot\mathbf R_A}V_{\mathrm{loc}}^A(\vG)\right]
  =-iG_xe^{-i\vG\cdot\mathbf R_A}V_{\mathrm{loc}}^A(\vG).
  \label{eq:gth-local-structure-derivative}
\end{equation}
$\vG=0$에서는 $G_x=0$이지만 constant term, charge convention과 Ewald compensation을
central hcore와 같은 방식으로 처리해야 한다.

AO matrix는
\begin{equation}
  V_{\mu\nu,\mathrm{loc}}^{\vk}
  =\langle\phi_{\mu\vk}|\hat V_{\mathrm{loc}}|\phi_{\nu\vk}\rangle
  \label{eq:gth-local-ao}
\end{equation}
이며 derivative에는 potential center $A$, bra AO center, ket AO center의 이동이 모두
포함된다. Translation identity를 사용하면 potential-center derivative를 AO derivative와
reciprocal structure-factor derivative의 조합으로 계산할 수 있다.

\subsection{Nonlocal projectors}

Nonlocal GTH operator는 angular momentum channel별 separable form
\begin{equation}
  \hat V_{\mathrm{nl}}^A
  =\sum_{l m}\sum_{ij}
   |p_{i,lm}^{A}\rangle h_{ij}^{A,l}
   \langle p_{j,lm}^{A}|
  \label{eq:gth-nonlocal-operator}
\end{equation}
이다. AO matrix element를 projector overlap
\begin{equation}
  P_{i,lm;\mu}^{A\vk}
  =\langle p_{i,lm}^{A}|\phi_{\mu\vk}\rangle
  \label{eq:gth-projector-overlap}
\end{equation}
로 쓰면
\begin{equation}
  V_{\mu\nu,\mathrm{nl}}^{\vk}
  =\sum_{A,lm,ij}
   P_{i,lm;\mu}^{A\vk *}h_{ij}^{A,l}
   P_{j,lm;\nu}^{A\vk}.
  \label{eq:gth-nonlocal-matrix}
\end{equation}
Product rule로
\begin{align}
  V_{\mu\nu,\mathrm{nl}}^{\vk,Ax}
  =\sum_{l m i j}
  \bigg[&P_{i,lm;\mu}^{A\vk,Ax *}h_{ij}^{A,l}
         P_{j,lm;\nu}^{A\vk}
  \notag\\
       +&P_{i,lm;\mu}^{A\vk *}h_{ij}^{A,l}
         P_{j,lm;\nu}^{A\vk,Ax}\bigg]
  +\text{AO-center terms}
  \label{eq:gth-nonlocal-derivative}
\end{align}
를 얻는다. Projector radial table $h_{ij}^{A,l}$ 자체는 원자와 함께 rigid하게 이동하므로
핵 좌표에 대한 parameter derivative는 없다.

\subsection{AO-row convention과 full matrix}

Integral engine은 흔히 AO row center에 대한 derivative matrix
$h_{\mu\nu}^{(1)}$를 반환한다. Atom $A$의 full matrix derivative를 만들려면
\begin{equation}
  h^{Ax}=h_{\mathrm{row},A}^{x}
       +h_{\mathrm{row},A}^{x\dagger}
       +h_{\mathrm{operator},A}^{x}
  \label{eq:hcore-full-derivative}
\end{equation}
처럼 bra/ket와 operator-center ownership을 합쳐야 한다. Complex k-point에서 두 번째
항은 단순 transpose가 아니라 conjugate transpose다.

Overlap derivative도 같은 AO-center logic으로
\begin{equation}
  S^{Ax}=S_{\mathrm{row},A}^{x}+S_{\mathrm{row},A}^{x\dagger}
  \label{eq:overlap-full-derivative}
\end{equation}
를 만족한다. $h^{Ax}$와 $S^{Ax}$는 explicit target gradient뿐 아니라 source HF CPHF
source 식~\eqref{eq:cphf-source}에도 사용된다.

\subsection{Pseudopotential 일관성}

Source density를 HF로 만들었다는 이유로 ``HF pseudopotential''이라는 별도의 analytic
형식이 생기지는 않는다. GTH-HF와 같은 label은 pseudopotential parameter set을 뜻한다.
Source HF와 target DFA energy/gradient는 동일한 GTH parameter와 valence ionic charge를
사용해야 한다. 서로 다른 pseudopotential로 source와 target을 계산하면 식
\eqref{eq:scope-dc-energy}의 한 energy surface를 미분하는 문제가 아니게 된다.

검증에서는 local, nonlocal, kinetic/overlap, ion--ion component를 분리해 FD와 비교한 뒤
전체 hcore derivative를 확인한다. 특히 nonlocal projector의 angular channel이 늘어나는
basis/pseudopotential에서 atom slice와 conjugation 오류가 잘 드러난다.

\section{FFTDF: reciprocal-space Coulomb와 exchange derivative}
\label{sec:fftdf}

\subsection{Uniform grid와 plane-wave representation}

FFTDF는 AO pair density를 uniform real-space grid에 평가하고 FFT를 통해 reciprocal
space로 옮긴다. k-point $\vk$의 density matrix가 만드는 total density Fourier
coefficient는 개략적으로
\begin{equation}
  \rho_D(\vG)
  =\frac{1}{N_k}\sum_{\vk\mu\nu}
   D_{\mu\nu}^{\vk}
   \int_\Omega e^{-i\vG\cdot\mathbf r}
   \phi_{\nu\vk}^{*}(\mathbf r)
   \phi_{\mu\vk}(\mathbf r)\dd\mathbf r.
  \label{eq:fftdf-density-fourier}
\end{equation}
Neutral periodic system의 Hartree energy는 $\vG=0$을 제외하고
\begin{equation}
  E_{\mathrm H}
  =\frac{1}{2\Omega}\sum_{\vG\ne0}
   \frac{4\pi}{G^2}|\rho_D(\vG)|^2
  \label{eq:fftdf-hartree-energy}
\end{equation}
로 표현된다. 역 FFT한 Hartree potential을 AO pair와 다시 적분하면 $J^{\vk}$를 얻는다.

\subsection{Nuclear derivative at fixed mesh}

Cell과 integer mesh를 고정한 nuclear derivative에서는 grid point와 $\vG$ set이 변하지
않는다. 따라서 Coulomb kernel $4\pi/G^2$의 nuclear derivative는 0이고 AO pair density만
미분된다.
\begin{equation}
  E_{\mathrm H}^{Ax}
  =\frac{1}{\Omega}\operatorname{Re}
   \sum_{\vG\ne0}\frac{4\pi}{G^2}
   \rho_D(\vG)^*\rho_D^{Ax}(\vG).
  \label{eq:fftdf-hartree-derivative}
\end{equation}
$\rho_D^{Ax}$에는 atom $A$에 속한 AO row와 column derivative, periodic image와 Bloch
phase convention이 들어간다. AO-row matrix를 먼저 만든 뒤 Hermitian reconstruction을
하면 식~\eqref{eq:hartree-gradient}의 full $J^{Ax}$를 얻는다.

\subsection{Exact exchange와 momentum transfer}

Source HF exchange는 pair density의 momentum transfer
\begin{equation}
  \vq=\vk'-\vk
  \label{eq:fftdf-momentum-transfer}
\end{equation}
마다 Coulomb kernel $4\pi/|\vG+\vq|^2$를 사용한다. Symbolically
\begin{equation}
  K_{\mu\nu}^{\vk}
  =\frac{1}{N_k}\sum_{\vk'}\sum_{\lambda\sigma}
   (\mu\vk,\lambda\vk'|\sigma\vk',\nu\vk)
   D_{\lambda\sigma}^{\vk'}.
  \label{eq:fftdf-exchange}
\end{equation}
Nuclear derivative는 네 AO center 중 이동하는 center에 product rule을 적용한다. Source
HF Hessian action에도 같은 FFT exchange provider를 transition density에 적용해야 한다.

\subsection{\texorpdfstring{$G=0$, $q=0$}{G=0, q=0}와 Ewald exchange divergence}

Hartree $\vG=0$ mode는 neutralizing convention과 GTH ionic charge에 맞추어 제외 또는
보정한다. Exchange에서는 $\vq=0$, $\vG=0$에서 Coulomb singularity가 나타나며 PySCF의
\texttt{exxdiv='ewald'} convention은 overlap-dependent correction을 더한다. 이 correction도
핵 이동에 따라 overlap이 변하므로
\begin{equation}
  K_{\mathrm{Ewald}}^{Ax}
  \propto S^{Ax}DS+SDS^{Ax}
  \label{eq:fftdf-ewald-exchange-derivative}
\end{equation}
형태의 derivative가 필요하다. Gamma라는 조건과 $\vq=0$이라는 조건을 혼동하면 sampled
k-point diagonal pair의 singular correction을 잃는다.

\subsection{Cutoff와 mesh}

Plane-wave kinetic cutoff는 대략
\begin{equation}
  E_{\mathrm{cut}}=\frac{G_{\max}^2}{2}
  \label{eq:cutoff-definition}
\end{equation}
이다. Orthorhombic cell에서 각 방향의 Nyquist wavevector를 고려하면 mesh size는 cell
길이에 비례하고 $\sqrt{E_{\mathrm{cut}}}$에 비례한다. PySCF는
\texttt{cell.ke\_cutoff}와 \texttt{cell.mesh}를 연동하지만 integer rounding 때문에
cutoff의 작은 변화가 mesh를
불연속적으로 바꿀 수 있다.

Full-energy FD에서 $\vR\pm h$는 반드시 같은 lattice와 같은 integer mesh를 사용해야
한다. Nuclear displacement마다 cutoff로 mesh를 다시 추정해 서로 다른 mesh를 사용하면
analytic derivative가 아닌 numerical representation 변화까지 FD에 섞인다. Metadata에는
cutoff와 resolved mesh를 둘 다 저장한다.

\subsection{Gradient 안에서 FFTDF가 소유하는 항}

FFTDF provider는 다음 quantity를 일관되게 제공해야 한다.

\begin{enumerate}
  \item source HF central $J[D]$와 $K[D]$,
  \item transition density에 대한 $\delta J$와 $\delta K$ for AX,
  \item fixed-density full $J^{Ax}[D]$와 $K^{Ax}[D]$,
  \item $\vq=0$ Ewald exchange-divergence derivative,
  \item target Hartree explicit derivative.
\end{enumerate}

Target XC는 NumInt/MGGA grid provider가 소유하고 GTH one-electron derivative는 별도
one-electron provider가 소유한다. 이 책임을 섞지 않아야 같은 mesh를 쓰면서도 component
double counting을 피할 수 있다.

\subsection{CPU와 GPU}

CPU는 PySCF periodic FFTDF central value와 project-local derivative matrix를 사용한다.
GPU는 GPU4PySCF FFT, AO grid와 CuPy contraction을 사용한다. GPU route에서 cell, lattice,
small pseudopotential table과 최종 reporting은 host에 남을 수 있지만 AO/grid FFT,
J/K, AX, Z-vector와 final electronic assembly가 device에 있어야 한다. Capability report는
이 host boundary를 명시하고 silent CPU fallback을 허용하지 않는다.

FFTDF validation은 $E_{\mathrm{cut}}$ 한 점에서 FD가 맞는 것만으로 끝나지 않는다.
Resolved mesh sequence에 대해 energy와 gradient가 수렴하고, Gamma와 equivalent one-k-point
representation이 일치하며, 식~\eqref{eq:translation-sum-rule}을 만족해야 한다.

\section{GDF: auxiliary metric과 nuclear response}
\label{sec:gdf}

\subsection{AO pair fitting}

Gaussian density fitting은 AO pair density를 atom-centered auxiliary function
$\{P\}$로 표현한다. k-point pair $(\vk_i,\vk_j)$의 momentum transfer를
$\vq=\vk_j-\vk_i$라 하면 three-center source와 two-center Coulomb metric은
\begin{align}
  B_{P,\mu\nu}^{ij}
  &=\left(P\vq\middle|\mu\vk_i,\nu\vk_j\right),
  \label{eq:gdf-three-center}\\
  M_{PQ}^{\vq}
  &=\left(P\vq\middle|Q\vq\right).
  \label{eq:gdf-two-center}
\end{align}
Metric factorization
\begin{equation}
  M^{\vq}=L^{\vq}L^{\vq\dagger}
  \label{eq:gdf-cholesky}
\end{equation}
을 사용하면 orthonormalized CDERI는
\begin{equation}
  C^{ij}=L^{-1}B^{ij}
  \label{eq:gdf-cderi}
\end{equation}
이다. Four-center integral은 CDERI 두 개의 contraction으로 근사된다.

\subsection{Three-center source derivative}

$B_{P,\mu\nu}$에는 auxiliary center, bra AO center, ket AO center가 있다. Atom $A$의
derivative는
\begin{align}
  B_{P,\mu\nu}^{ij,Ax}
  ={}&\delta_{A,A_P}\partial_{P_x}B_{P,\mu\nu}^{ij}
    +\delta_{A,A_\mu}\partial_{\mu_x}B_{P,\mu\nu}^{ij}
  \notag\\
   &+\delta_{A,A_\nu}\partial_{\nu_x}B_{P,\mu\nu}^{ij}
  \label{eq:gdf-source-centers}
\end{align}
이다. Integral library의 \texttt{ip1}/\texttt{ip2}가 electronic-coordinate derivative를 반환하면
nuclear-center derivative 부호는 translational identity에 따라 반대가 될 수 있다.
각 axis와 atom slice를 작은 component test로 확인해야 한다.

Periodic GDF builder는 short-range real-space part, long-range reciprocal-space part와
$\vq=0$ correction을 합칠 수 있다.
\begin{equation}
  B^{ij}=B_{\mathrm{SR}}^{ij}+B_{\mathrm{LR}}^{ij}
  +\delta_{\vq,0}B_{vbar}^{ij}.
  \label{eq:gdf-source-parts}
\end{equation}
Derivative도 central builder와 동일한 screening, lattice image와 reciprocal mesh를
사용해야 한다.

\subsection{Metric과 Cholesky derivative}

Two-center metric derivative에는 두 auxiliary center response가 있다.
\begin{equation}
  M_{PQ}^{\vq,Ax}
  =\delta_{A,A_P}\partial_{P_x}M_{PQ}^{\vq}
   +\delta_{A,A_Q}\partial_{Q_x}M_{PQ}^{\vq}.
  \label{eq:gdf-metric-centers}
\end{equation}
식~\eqref{eq:gdf-cholesky}을 미분하면
\begin{equation}
  M^{Ax}=L^{Ax}L^\dagger+LL^{Ax\dagger}.
  \label{eq:gdf-cholesky-derivative}
\end{equation}
$W=L^{-1}M^{Ax}L^{-\dagger}$를 정의하고 lower-Cholesky gauge에서
\begin{equation}
  L^{Ax}=L\,\Phi(W),
  \label{eq:gdf-cholesky-gauge}
\end{equation}
로 쓴다. $\Phi$는 strict lower triangle과 diagonal의 절반을 취한다. 따라서 CDERI
derivative는
\begin{equation}
  C^{Ax}
  =L^{-1}B^{Ax}-L^{-1}L^{Ax}C
  \label{eq:gdf-cderi-derivative}
\end{equation}
이다. 두 번째 항이 metric/auxiliary response다. 이를 생략하면 central fitted energy와
다른 functional form을 미분하게 된다.

\subsection{J와 K derivative}

Diagonal k-pair CDERI로 Coulomb auxiliary density를
\begin{equation}
  j_P=\frac{1}{N_k}\sum_{\vk\mu\nu}
      C_{P,\mu\nu}^{\vk\vk}D_{\nu\mu}^{\vk}
  \label{eq:gdf-coulomb-aux-density}
\end{equation}
라 하면 fixed-density derivative는
\begin{equation}
  j_P^{Ax}=\frac{1}{N_k}\sum_{\vk\mu\nu}
      C_{P,\mu\nu}^{\vk\vk,Ax}D_{\nu\mu}^{\vk}.
  \label{eq:gdf-coulomb-aux-derivative}
\end{equation}
이에 따라
\begin{equation}
  J_{\mu\nu}^{\vk,Ax}
  =\sum_P\left(j_P^{Ax}C_{P,\mu\nu}^{\vk\vk}
              +j_PC_{P,\mu\nu}^{\vk\vk,Ax}\right).
  \label{eq:gdf-j-derivative}
\end{equation}
Exchange도 두 CDERI factor 각각에 product rule을 적용하며 canonical k-pair의 reverse
block은 conjugate/transposed relation으로 복원한다.

\subsection{Auxiliary response가 gradient에서 나타나는 두 위치}

첫째, target Hartree explicit derivative 식~\eqref{eq:hartree-gradient}에
$C^{Ax}$가 들어간다. 둘째, source HF CPHF source와 Hessian action에 GDF $J/K$
derivative가 들어간다. Target explicit에서만 metric response를 더하고 source HF
response에서 빠뜨리거나 그 반대로 하면 Z-vector correction이 full-energy FD와 맞지
않는다.

\subsection{Value parity guard}

Derivative builder가 재구성한 central CDERI와 실제 SCF provider가 저장한 CDERI의 차이를
\begin{equation}
  \Delta_C=\max|C_{\mathrm{response\ builder}}-C_{\mathrm{SCF\ storage}}|
  \label{eq:gdf-cderi-parity}
\end{equation}
로 정의한다. $\Delta_C$가 tolerance를 넘으면 $B^{Ax}$와 $M^{Ax}$가 수학적으로
그럴듯해도 실제 energy representation의 derivative가 아니므로 gradient를 거부한다.
원인은 흔히 다른 $\omega$, reciprocal mesh, auxiliary basis ordering, k-pair orientation,
metric factorization gauge다.

\subsection{CPU/GPU 경계}

CPU package는 PySCF periodic integral engine으로 central/derivative primitive를 만들고
NumPy/SciPy linear algebra를 사용한다. GPU package는 periodic 2-center/3-center integral
primitive 생성은 host boundary로 명시하지만, CDERI storage, Cholesky/metric response,
$dJ/dR$, $dK/dR$, AX/Z-vector와 final contraction은 CuPy device에서 수행한다.
\texttt{device="gpu"}는 integral metadata까지 모두 GPU라는 뜻이 아니라, 지원한다고 선언한
numerical contraction이 CPU implementation으로 조용히 우회하지 않는다는 뜻이다.

GDF completion은 H$_2$뿐 아니라 distorted H$_2$O, Gamma와 non-Gamma k-mesh, SZV와
DZVP 수준에서 full-energy FD, translation sum, CDERI parity를 함께 통과해야 한다.

\section{SCAN/r2SCAN meta-GGA energy와 gradient}
\label{sec:mgga}

\subsection{Semilocal ingredients}

Spin-unpolarized MGGA는 한 grid point에서
\begin{equation}
  e_{\mathrm{xc}}
  =f(\rho,\sigma,\tau),
  \qquad \sigma=\nabla\rho\cdot\nabla\rho
  \label{eq:mgga-ingredients}
\end{equation}
를 평가한다. Variational derivative에 필요한 weights를
\begin{equation}
  v_\rho=\frac{\partial f}{\partial\rho},
  \qquad
  v_\sigma=\frac{\partial f}{\partial\sigma},
  \qquad
  v_\tau=\frac{\partial f}{\partial\tau}
  \label{eq:mgga-weights}
\end{equation}
로 둔다.

SCAN은 균일 전자 기체, 단일 orbital, 완만히 변하는 density의 limit를 잇도록 설계된
constraint-based functional이다. r2SCAN은 SCAN switching function의 수치적 민감도를
완화해 grid convergence를 개선하지만 $\rho$, $\sigma$, $\tau$를 사용하는 derivative
structure는 같다.

\subsection{AO representation of \texorpdfstring{$\tau$}{tau}}

식~\eqref{eq:tau-definition}을 AO product로 쓰면
\begin{equation}
  \tau(\mathbf r)
  =\frac{1}{2N_k}\sum_{\vk\mu\nu}D_{\mu\nu}^{\vk}
   \sum_{y=x,y,z}
   \partial_y\phi_{\nu\vk}^{*}(\mathbf r)
   \partial_y\phi_{\mu\vk}(\mathbf r).
  \label{eq:mgga-tau-ao}
\end{equation}
Target MGGA Fock의 tau contribution은
\begin{equation}
  V_{\mu\nu}^{\tau,\vk}
  =\frac12\int_\Omega v_\tau(\mathbf r)
   \nabla\phi_{\nu\vk}^{*}(\mathbf r)\cdot
   \nabla\phi_{\mu\vk}(\mathbf r)\dd\mathbf r.
  \label{eq:mgga-vtau-matrix}
\end{equation}
이다. $1/2$는 식~\eqref{eq:tau-definition}에서 온다. Numerical integration library가
이미 $\tau$ convention을 반영한 effective weight를 반환하는지 확인한 뒤 factor를 한
번만 적용한다.

\subsection{Explicit nuclear derivative}

Density matrix를 고정하면
\begin{equation}
  E_{\mathrm{xc}}^{Ax}
  =\int_\Omega\left(
    v_\rho\rho^{Ax}
    +v_\sigma\sigma^{Ax}
    +v_\tau\tau^{Ax}\right)\dd\mathbf r
  \label{eq:mgga-xc-derivative}
\end{equation}
이다. Fixed-cell uniform grid에서는 quadrature point와 volume weight가 핵 좌표에 따라
변하지 않으므로 AO derivative만 남는다. GGA 부분은
\begin{equation}
  \sigma^{Ax}=2\nabla\rho\cdot\nabla\rho^{Ax}
  \label{eq:mgga-sigma-derivative}
\end{equation}
이고 tau derivative는
\begin{align}
  \tau^{Ax}(\mathbf r)
  =\frac{1}{2N_k}\sum_{\vk\mu\nu}D_{\mu\nu}^{\vk}
   \sum_y\big[&
   \partial_y\phi_{\nu\vk}^{Ax *}
   \partial_y\phi_{\mu\vk}
  \notag\\
   &+\partial_y\phi_{\nu\vk}^{*}
   \partial_y\phi_{\mu\vk}^{Ax}\big].
  \label{eq:mgga-tau-derivative}
\end{align}
AO center derivative와 spatial gradient가 함께 있으므로 second spatial derivative AO가
필요하다. 이 점이 GGA보다 MGGA gradient kernel이 복잡한 주된 이유다.

\subsection{Nonvariational response와의 결합}

MGGA에서도 source orbital은 HF가 결정한다. 따라서 Hessian $A$는 식
\eqref{eq:hf-hessian-action}과 동일하고, target Fock에 식
\eqref{eq:mgga-vtau-matrix}가 더해져 target VO gradient
\begin{equation}
  g_{ai}^{\vk,\mathrm{MGGA}}
  =\left[C_v^{\vk\dagger}
   \left(h+J+V_\rho+V_\sigma+V_\tau\right)^{\vk}
   C_o^{\vk}\right]_{ai}
  \label{eq:mgga-target-vo}
\end{equation}
를 바꾼다. Z-vector solve 후 식~\eqref{eq:response-gradient}에 source HF nuclear
perturbation을 contraction한다. MGGA kernel을 explicit XC term에만 연결하고 target Fock
VO RHS를 PBE 값으로 남기면 nonvariational response가 틀린다.

\subsection{Grid와 aliasing}

Gaussian AO의 high-frequency component는 grid가 거칠면 aliasing된다. $\tau$는 AO
gradient product이므로 density보다 높은 reciprocal component에 민감하다. 따라서
MGGA에서 \texttt{cell.mesh} 또는 effective integration mesh convergence를 별도로 확인해야
한다. Mesh를 한 방향씩 바꾸지 말고 resolved integer triplet과 cutoff를 함께 기록한다.

Translation sum rule은 grid aliasing을 민감하게 보여준다. Analytic--FD component error가
작아도
\begin{equation}
  \left\|\sum_A\mathbf g_A^{\mathrm{xc}}\right\|
  \label{eq:mgga-xc-translation}
\end{equation}
가 크면 explicit XC grid derivative가 아직 수렴하지 않은 것이다.

\subsection{CPU/GPU parity}

CPU와 GPU는 같은 libxc functional label, density convention, mesh, AO derivative order를
사용해야 한다. 비교 항목은 total energy와 final gradient뿐 아니라
$E_{\mathrm{xc}}$, target Fock VO block, $v_\tau$ row derivative, Z-vector RHS와 residual이다.
SCAN/r2SCAN completion case에는 Gamma H$_2$, distorted H$_2$O와 sampled k-point case가
포함되어야 한다. 또한 target MGGA와 source-HF density-fitting provider는 서로 다른
책임을 가지므로 FFTDF와 GDF만이 아니라 RSDF에서도 같은 $\rho$, $\nabla\rho$, $\tau$
target path를 검증해야 한다. RSDF+MGGA case는 Eq.~\eqref{eq:mgga-xc-derivative}의
XC derivative와 Eq.~\eqref{eq:rsdf-total-gradient}의 source/metric auxiliary response가
동시에 작동하는 composition test이다.

\section{RSDF: range-separated source와 metric response}
\label{sec:rsdf}

\subsection{Range separation의 목적}

Range-separated density fitting(RSDF)은 Coulomb kernel의 exact identity
\begin{equation}
  \frac1r
  =\frac{\operatorname{erfc}(\omega r)}{r}
   +\frac{\operatorname{erf}(\omega r)}{r}
  \equiv v_{\mathrm{SR}}(r;\omega)+v_{\mathrm{LR}}(r;\omega)
  \label{eq:rsdf-range-separation}
\end{equation}
를 사용한다. Short-range part는 real-space lattice sum에서 빠르게 감쇠하고 long-range
part는 smooth reciprocal-space sum으로 계산하기 쉽다. Infinite sum에서는 $\omega$와
무관하지만 finite screening과 finite mesh에서는 $\omega$가 numerical representation의
일부다.

\subsection{Source와 metric parameter는 독립적이다}

Three-center source builder와 two-center metric builder가 최적화하는 exponent 분포가
다르므로 PySCF RSDF는 서로 다른 parameter를 선택할 수 있다. 이를
\begin{align}
  B^{ij}
  &=B_{\mathrm{SR}}^{ij}(\omega_B,\epsilon_B)
    +B_{\mathrm{LR}}^{ij}(\omega_B,\mathcal G_B)
    +\delta_{\vq,0}B_{vbar}^{ij},
  \label{eq:rsdf-source-split}\\
  M^{\vq}
  &=M_{\mathrm{SR}}^{\vq}(\omega_M,\epsilon_M)
    +M_{\mathrm{LR}}^{\vq}(\omega_M,\mathcal G_M)
    +\delta_{\vq,0}M_0
  \label{eq:rsdf-metric-split}
\end{align}
로 쓴다. 일반적으로
\begin{equation}
  \omega_B\ne\omega_M,
  \qquad \mathcal G_B\ne\mathcal G_M.
  \label{eq:rsdf-independent-parameters}
\end{equation}
이다. Source parameter를 metric derivative에 재사용하는 것은 approximation이 아니라
다른 central CDERI를 미분하는 오류다.

\subsection{Fixed-cell nuclear derivative}

Cell, resolved source/metric mesh와 $\omega_B$, $\omega_M$를 고정하면
\begin{align}
  B^{ij,Ax}
  &=B_{\mathrm{SR}}^{ij,Ax}
    +B_{\mathrm{LR}}^{ij,Ax}
    +\delta_{\vq,0}B_{vbar}^{ij,Ax},
  \label{eq:rsdf-source-derivative}\\
  M^{\vq,Ax}
  &=M_{\mathrm{SR}}^{\vq,Ax}
    +M_{\mathrm{LR}}^{\vq,Ax}
    +\delta_{\vq,0}M_0^{Ax}.
  \label{eq:rsdf-metric-derivative}
\end{align}
Short-range derivative는 central source/metric과 같은 lattice image list 및 screening
precision을 사용한다. Long-range derivative는 같은 reciprocal vectors와 structure
factor를 사용한다. Nuclear derivative에서 integer mesh는 변하지 않는다.

Resolved metric을 Cholesky factorization한 뒤 식~\eqref{eq:gdf-cderi-derivative}을
그대로 적용한다. 즉 RSDF에서도
\begin{equation}
  C^{Ax}=L^{-1}B^{Ax}-L^{-1}L^{Ax}C
  \label{eq:rsdf-cderi-derivative}
\end{equation}
이며 source와 metric response를 모두 보존해야 한다.

\subsection{\texorpdfstring{$q=0$}{q=0} correction}

$\vq=0$은 Gamma-only를 뜻하지 않는다. Sampled k-point의 $(\vk_i,\vk_i)$ diagonal pair도
$\vq=0$이다. 따라서
\begin{equation}
  \delta_{\vq,0}
  =\begin{cases}
    1,&\|\vk_j-\vk_i-\vG\|<\epsilon_q\ \text{for some reciprocal }\vG,\\
    0,&\text{otherwise}
   \end{cases}
  \label{eq:rsdf-qzero-test}
\end{equation}
처럼 reciprocal equivalence를 고려한다. Source $vbar$, metric $M_0$, Ewald exchange
correction의 value와 derivative가 같은 pair orientation을 사용해야 한다.

\subsection{Canonical pair storage}

Memory를 줄이기 위해 canonical $(i,j)$ pair만 저장할 수 있다. Reverse pair는
\begin{equation}
  C_{P,\nu\mu}^{ji}
  =C_{P,\mu\nu}^{ij *}
  \label{eq:rsdf-reverse-pair}
\end{equation}
로 복원하고 derivative도
\begin{equation}
  C_{P,\nu\mu}^{ji,Ax}
  =C_{P,\mu\nu}^{ij,Ax *}
  \label{eq:rsdf-reverse-pair-derivative}
\end{equation}
를 따라야 한다. AO transpose만 하고 complex conjugate를 빠뜨리면 Gamma test는 통과할
수 있지만 non-Gamma k-point에서 imaginary leakage와 gradient error가 발생한다.

\subsection{RSDF response가 사용되는 곳}

RSDF CDERI derivative는 GDF와 마찬가지로 target Hartree explicit derivative, source HF
$J/K$ nuclear derivative와 CPHF source에 사용된다. AX operation은 central RSDF J/K를
transition density에 적용한다. 전체 구조는
\begin{equation}
  g_{Ax}^{\mathrm{RSDF}}
  =g_{Ax}^{1e}+g_{Ax}^{\mathrm H}[C,C^{Ax}]
   +g_{Ax}^{\mathrm{xc}}+g_{Ax}^{\mathrm{Pulay,t}}
   +g_{Ax}^{\mathrm{NN}}+z^\Herm B_{\mathrm{HF}}^{Ax}[C,C^{Ax}]
  \label{eq:rsdf-total-gradient}
\end{equation}
로 요약할 수 있다.

\subsection{GPU host boundary}

현재 GPU4PySCF에는 RSDF periodic two-/three-center derivative primitive를 모두 만드는
device integral engine이 없다. 따라서 host에서 central/derivative integral metadata를
만들고 CuPy로 옮긴 뒤 metric factorization, $dC$, $dJ/dR$, $dK/dR$, AX/Z-vector와 final
assembly를 device에서 수행한다. 이 경계는 metadata에 명시하며 CPU DC-DFT gradient
kernel로의 silent fallback과 구분한다.

\subsection{RSDF completion checks}

RSDF validation은 다음 순서를 따른다.

\begin{enumerate}
  \item 실제 provider가 resolve한 $\omega_B$, source mesh, $\omega_M$, metric mesh를 기록한다.
  \item Response builder의 central CDERI와 SCF storage의 value parity를 확인한다.
  \item Source-only, metric-only, total $dC$를 component FD와 비교한다.
  \item $dJ/dR$, $dK/dR$, target explicit와 Z response를 각각 검사한다.
  \item Gamma 및 sampled k-point full-energy central FD를 수행한다.
  \item CPU/GPU parity와 translation sum rule을 확인한다.
\end{enumerate}

Auxiliary basis 또는 source/metric mesh를 바꿔도 error가 특정 molecule에서만 작아지는
경험적 correction은 completion evidence로 인정하지 않는다.

\section{수치 검증, convergence와 보고 단위}
\label{sec:validation}

\subsection{Full-energy central finite difference}

Analytic gradient의 authoritative numerical reference는
\begin{equation}
  g_{Ax}^{\mathrm{FD}}(h)
  =\frac{E_{\mathrm{DC-DFA}}(\vR+h\mathbf e_{Ax})
         -E_{\mathrm{DC-DFA}}(\vR-h\mathbf e_{Ax})}{2h}
  \label{eq:central-fd}
\end{equation}
이다. 두 displaced geometry 각각에서 source HF를 같은 convergence tolerance로 다시
수렴시키고 target DFA energy를 새로 평가한다. Central density 또는 MO를 그대로 두고
operator만 미분한 값은 component diagnostic이지 식~\eqref{eq:central-fd}가 아니다.

Leading truncation error는 $O(h^2)$이고 너무 작은 $h$에서는 SCF와 floating-point noise가
$O(\epsilon/h)$로 증폭된다. 여러 step $h$, $h/2$, $h/4$에서 plateau를 찾고 analytic
gradient가 그 plateau와 일치하는지 본다.

\subsection{Error norm}

Selected component 집합 $\mathcal I$에 대해
\begin{align}
  \Delta_{\max}
  &=\max_{(A,x)\in\mathcal I}|g_{Ax}-g_{Ax}^{\mathrm{FD}}|,
  \label{eq:fd-max-error}\\
  \Delta_{\mathrm{RMS}}
  &=\sqrt{\frac1{|\mathcal I|}\sum_{(A,x)\in\mathcal I}
          |g_{Ax}-g_{Ax}^{\mathrm{FD}}|^2}
  \label{eq:fd-rms-error}
\end{align}
를 기록한다. 한 축 또는 한 원자만 맞는 결과를 full-gradient validation이라고 부르지
않는다. Completion gate는 distorted geometry의 모든 $3N_{\mathrm{atom}}$ component를
포함한다.

\subsection{SCF와 Z-vector residual}

Displaced HF의 energy/density tolerance가 gradient target보다 충분히 작아야 한다. HF
orbital residual, density change와 energy change를 함께 기록한다. Z-vector는 식
\eqref{eq:zvector-residual}의 maximum residual과 condition estimate를 저장한다. CPU와 GPU
solver가 같은 gradient를 내더라도 서로 다른 큰 residual에 cancellation으로 도달했다면
parity 근거가 아니다.

\subsection{Representation convergence}

Analytic--FD agreement를 확인한 다음 다음 numerical parameter를 한 번에 하나씩
수렴시킨다.

\begin{itemize}
  \item primary Gaussian basis와 GTH pseudopotential,
  \item FFTDF cutoff와 resolved integer mesh,
  \item target XC integration mesh,
  \item GDF/RSDF auxiliary basis,
  \item GDF screening와 reciprocal mesh,
  \item RSDF source $\omega_B$/mesh와 metric $\omega_M$/mesh,
  \item k-point mesh.
\end{itemize}

Cutoff를 보고할 때 mesh를 생략하지 않는다. 같은 cutoff라도 lattice 길이가 달라지면
integer mesh가 달라진다. RSDF는 source와 metric parameter를 하나의
\texttt{omega} 또는 \texttt{mesh}로 축약해 저장하지 않는다.

\subsection{Symmetry-independent checks}

Translation sum rule error를
\begin{equation}
  \Delta_{\mathrm{trans}}
  =\left\|\sum_A\mathbf g_A\right\|_2
  \label{eq:translation-error}
\end{equation}
로 정의한다. Symmetric H$_2$만 사용하면 누락 항이 대칭으로 상쇄될 수 있으므로 bond,
angle과 전체 orientation을 조금 왜곡한 H$_2$O를 필수로 사용한다. Gamma와 sampled
k-point, 실수와 complex tensor 경로를 모두 포함한다.

Equivalent Gamma representation은
\begin{equation}
  E_{\Gamma}=E_{N_k=1},
  \qquad
  \mathbf g_{\Gamma}=\mathbf g_{N_k=1}
  \label{eq:gamma-one-kpoint-parity}
\end{equation}
를 만족해야 한다. 이 검사는 singleton k-axis를 제거하거나 유지하는 shape adapter의
오류를 찾는다.

\subsection{CPU/GPU parity}

같은 cell, basis, pseudopotential, mesh, k-point와 SCF tolerance에서
\begin{align}
  \Delta_E^{\mathrm{CPU/GPU}}
  &=|E_{\mathrm{CPU}}-E_{\mathrm{GPU}}|,
  \label{eq:cpu-gpu-energy-error}\\
  \Delta_g^{\mathrm{CPU/GPU}}
  &=\max_{Ax}|g_{Ax}^{\mathrm{CPU}}-g_{Ax}^{\mathrm{GPU}}|
  \label{eq:cpu-gpu-gradient-error}
\end{align}
를 기록한다. GDF/RSDF host integral generation은 명시적인 GPU host boundary일 수 있지만
GPU path가 CPU package의 final gradient를 호출해서 같은 값을 반환하는 것은 parity가
아니라 fallback이다. Capability와 timing metadata로 ownership을 확인한다.

\subsection{초기 acceptance threshold}

이 project의 초기 completion target은 다음과 같다.

\begin{center}
\begin{tabular}{ll}
\toprule
Test & Threshold \\
\midrule
Analytic vs converged full-energy FD & $\Delta_{\max}\le10^{-5}$ Ha/Bohr \\
CPU/GPU energy & $\Delta_E\le10^{-8}$ Ha \\
CPU/GPU gradient & $\Delta_g\le10^{-7}$ Ha/Bohr \\
Translation norm & $\Delta_{\mathrm{trans}}\le10^{-6}$ Ha/Bohr \\
\bottomrule
\end{tabular}
\end{center}

Threshold는 basis/cutoff/FD-step convergence evidence와 함께 해석한다. Coarse mesh에서
analytic과 FD가 정확히 일치해도 그 gradient가 converged physical gradient라는 뜻은
아니다.

\subsection{Result metadata}

Energy/gradient result는 최소한 다음 정보를 machine-readable form으로 보존한다.

\begin{itemize}
  \item method, target XC, source reference, charge/spin,
  \item atom coordinates와 lattice 및 각 unit,
  \item basis와 pseudopotential,
  \item Gamma/k-point list와 weights,
  \item DF backend와 auxiliary basis,
  \item cutoff, resolved FFT/XC mesh,
  \item RSDF source/metric settings,
  \item HF convergence와 energy timings,
  \item Z dimension/residual/condition과 gradient component timings,
  \item device, host boundaries와 silent-fallback flag,
  \item gradient array unit와 force-sign convention.
\end{itemize}

Canonical gradient output은 NumPy array in Ha/Bohr다. JSON에는 큰 array를 중복 저장하기보다
shape, dtype, unit, checksum과 NPY path를 기록할 수 있다.

\subsection{흔한 failure signature}

\begin{center}
\begin{longtable}{p{0.35\textwidth}p{0.55\textwidth}}
\toprule
관찰 & 우선 확인할 원인 \\
\midrule
Gamma는 맞고 k-point가 틀림 & $1/N_k$, conjugation, reverse pair, $q=0$ \\
작은 system은 맞고 큰 system에서 GDF 오차 증가 & metric/auxiliary-center response \\
Cutoff에 따라 error가 불연속 & displaced geometry의 resolved mesh 불일치 \\
MGGA만 net translation이 큼 & $\tau$ factor, second AO derivative, grid aliasing \\
RSDF value parity 실패 & source/metric $\omega$ 또는 mesh 혼용 \\
Z residual은 작고 전체 FD만 틀림 & explicit hcore/J/XC/Pulay derivative \\
CPU/GPU total만 같고 component가 다름 & silent fallback 또는 서로 다른 ownership \\
\bottomrule
\end{longtable}
\end{center}

\section{Target multigrid와 source DF의 독립적 소유권}
\label{sec:target-multigrid}

\subsection{두 개의 선택 축}

Periodic DC-DFT 계산에는 서로 다른 역할을 하는 두 numerical 선택이 있다. 첫 번째는
source HF의 Coulomb/exchange와 orbital Hessian action을 계산하는
\textbf{source two-electron backend}다. FFTDF, GDF, RSDF 중 하나가 이 역할을 맡는다.
두 번째는 수렴된 HF density에서 target DFA의 $J$와 $E_{\mathrm{xc}}$를 평가하는
\textbf{target integrator}다. 두 선택을 하나의 ``DF method''로 합쳐 부르면 source와
target의 책임이 뒤섞인다.

이를 다음 ordered pair로 적는다.
\begin{equation}
  \mathcal R
  =\left(\mathcal B_{\mathrm{source}},\mathcal I_{\mathrm{target}}\right),
  \qquad
  \mathcal B_{\mathrm{source}}\in\{\mathrm{FFTDF},\mathrm{GDF},\mathrm{RSDF}\}.
  \label{eq:target-owner-split}
\end{equation}
Target integrator를 바꾸어도 source HF density $D^{\mathrm{HF}}$, HF energy, occupied/virtual
orbital과 source Hessian $A^{\mathrm{HF}}$는 변하지 않아야 한다. 바뀌는 것은
$E_{\mathrm{DFA}}[D^{\mathrm{HF}}]$의 target quadrature와 그로부터 만들어지는 target
Fock, orbital-gradient right-hand side 및 explicit target derivative다.

\subsection{세 target integration mode}

현재 GPU route는 다음 세 mode를 구분한다.

\begin{center}
\begin{longtable}{p{0.21\textwidth}p{0.22\textwidth}p{0.22\textwidth}p{0.23\textwidth}}
\toprule
Mode & Target $J$ & Target XC & Source HF/AX \\
\midrule
\texttt{uniform} & selected DF & uniform KNumInt & selected DF \\
\texttt{multigrid\_xc} & selected DF & multigrid & selected DF \\
\texttt{multigrid\_j\_xc} & multigrid & multigrid & selected DF \\
\bottomrule
\end{longtable}
\end{center}

\texttt{multigrid\_xc}를 M1이라 부르면 target effective potential은
\begin{equation}
  V_{\mathrm{eff}}^{\mathrm{M1}}[D]
  =J_{\mathcal B_{\mathrm{source}}}[D]
   +V_{\mathrm{xc}}^{\mathrm{MG}}[D]
  \label{eq:target-m1-veff}
\end{equation}
이다. 따라서 finite auxiliary basis를 사용하는 GDF/RSDF target $J$의 representation을
유지하면서 XC quadrature만 multigrid로 교체한다.

반면 \texttt{multigrid\_j\_xc}를 M2라 부르면
\begin{equation}
  V_{\mathrm{eff}}^{\mathrm{M2}}[D]
  =J_{\mathrm{MG}}[D]+V_{\mathrm{xc}}^{\mathrm{MG}}[D].
  \label{eq:target-m2-veff}
\end{equation}
M2의 $J_{\mathrm{MG}}$와 finite-auxiliary GDF/RSDF $J$는 서로 다른 수치
representation이다. Mesh와 auxiliary basis를 각각 수렴시키면 같은 Coulomb limit로
가지만, finite setting에서 energy가 같아야 할 이유는 없다. 따라서 M2는 단순한
GDF/RSDF drop-in acceleration이 아니라 별도 target model로 보고 비교해야 한다.

\subsection{Energy와 gradient ownership}

Mode에 따라 analytic gradient에서 target Coulomb derivative를 정확히 한 번만 더해야
한다. Ownership predicate를 $\chi_J(\mathcal I)$라 하면
\begin{equation}
  \mathbf g_{\mathrm{target}}
  =\begin{cases}
    \mathbf g_J^{\mathcal B_{\mathrm{source}}}
    +\mathbf g_{\mathrm{xc}}^{\mathrm{uniform}},
      &\mathcal I=\texttt{uniform},\\
    \mathbf g_J^{\mathcal B_{\mathrm{source}}}
    +\mathbf g_{\mathrm{xc}}^{\mathrm{MG}},
      &\mathcal I=\texttt{multigrid\_xc},\\
    \mathbf g_{J+\mathrm{xc}}^{\mathrm{MG}},
      &\mathcal I=\texttt{multigrid\_j\_xc}.
  \end{cases}
  \label{eq:target-gradient-owner}
\end{equation}
마지막 줄에서 selected-DF target $dJ/dR$를 다시 더하면 Coulomb response가 이중 계산된다.
반대로 M1에서 selected-DF $dJ/dR$를 생략하면 energy와 gradient가 서로 다른 Hamiltonian을
미분하게 된다.

Target hcore와 GTH pseudopotential은 세 mode 모두 selected one-electron provider가
소유한다. Multigrid object를 target KS object 전체에 암묵적으로 부착하면 upstream
\texttt{get\_hcore()} 또는 \texttt{get\_j()} dispatch까지 바뀔 수 있다. 따라서 구현은
low-level target XC/J 결과를 명시적으로 조합하고, hcore/PP ownership을 별도 metadata로
보존해야 한다.

\subsection{왜 target multigrid가 HF Hessian을 바꾸지 않는가}

DC-DFT의 adjoint equation은 source HF stationarity constraint를 미분한 것이다.
그러므로
\begin{equation}
  A_{\mathrm{HF}}^Tz=-g_{\mathrm{target}}
\end{equation}
에서 오른쪽 $g_{\mathrm{target}}$은 target integrator에 따라 달라지지만 왼쪽
$A_{\mathrm{HF}}$는 source backend를 유지한다. M1/M2를 선택했다고 해서 AX iteration의
$J/K$를 multigrid로 바꾸면, 더 이상 원래 HF density가 만족한 stationarity constraint의
response를 푸는 것이 아니다.

이 구분은 계산 흐름을 다음처럼 만든다.
\begin{enumerate}
  \item GDF 또는 RSDF로 source HF를 수렴시켜 $D^{\mathrm{HF}}$를 만든다.
  \item 선택한 target integrator로 $E_{\mathrm{DFA}}[D^{\mathrm{HF}}]$와
        $F_{\mathrm{target}}$을 만든다.
  \item $F_{\mathrm{target},vo}$에서 $g_{\mathrm{target}}$을 만든다.
  \item source GDF/RSDF AX로 $z$를 푼다.
  \item target explicit derivative와 source-HF response derivative를 조립한다.
\end{enumerate}

\subsection{Source mesh와 target mesh}

\texttt{cell.ke\_cutoff}로 resolve한 \texttt{cell.mesh}는 FFT 기반 source operator,
GTH local potential 및 여러 reciprocal-space 작업에 사용될 수 있다. Target multigrid
mesh는 target XC/J quadrature를 위한 별도 numerical parameter다. 두 mesh가 같은
기본값을 가질 수는 있지만 같은 정확도 의미를 갖지는 않는다.

Source와 target convergence를 분리해
\begin{equation}
  \lim_{\mathcal M_s\to\infty}
  \lim_{\mathcal M_t\to\infty}
  \mathbf g(\mathcal M_s,\mathcal M_t)
  \label{eq:target-grid-convergence}
\end{equation}
를 확인한다. 하나의 cutoff만 높이고 target mesh를 보고하지 않으면 어느
representation이 수렴한 것인지 알 수 없다. \texttt{target\_multigrid\_mesh=null}은
자동 memory cap을 뜻하지 않고 현재 cell mesh를 그대로 상속한다. 큰 cell에서는 resolved
triplet과 grid-point 수를 반드시 report해야 한다.

Multigrid는 grid를 여러 Gaussian/locality level로 나누지만 finite quadrature라는 사실은
변하지 않는다. 따라서 촘촘한 uniform grid보다 항상 정확하거나 translation invariant한
방법이라고 가정할 수 없다. Energy phase scan, translation sum과 full-gradient mesh
convergence를 각각 확인한다.

\subsection{Energy workspace 재사용과 수명}

Target energy에서 만든 multigrid task, sorted Gaussian pair와 target Fock은 gradient에서
재사용할 수 있다. Energy와 gradient가 서로 다른 multigrid object 또는 mesh를 사용하면
central energy의 derivative가 아니다. 반대로 target XC derivative가 끝난 뒤에도 큰 task
cache를 계속 보유하면 뒤따르는 GDF/RSDF response와 memory peak가 겹친다.

안전한 lifetime은 다음과 같다.
\begin{equation}
  \text{build at target energy}
  \longrightarrow
  \text{reuse at target gradient}
  \longrightarrow
  \text{release before DF response}.
  \label{eq:target-multigrid-memory}
\end{equation}
이는 grid point를 버리는 근사가 아니라, 마지막 consumer가 끝난 object의 소유권을
종료하는 memory 관리다.

\subsection{권장 사용 범위}

M1은 source HF와 finite-auxiliary target $J$를 유지한 채 target XC만 교체하므로 현재
가장 보수적인 opt-in route다. M2는 analytic하게 일관될 수 있지만 target $J$
representation도 바꾸므로 별도 convergence가 필요하다. PP-only multigrid는 central
hcore, explicit PP derivative, nuclear CPHF source 및 Z contraction을 모두 같은
representation으로 옮기기 전에는 사용할 수 없다. 일부 PP term만 바꾸면 energy와
gradient ownership이 불일치한다.

\section{Matrix-free Z-vector와 bounded-memory response}
\label{sec:matrix-free-bounded-memory}

\subsection{Response dimension과 dense matrix 비용}

k-point $\vk$에서 occupied--virtual rotation의 complex dimension은
$n_o^{\vk}n_v^{\vk}$다. 식~\eqref{eq:realified-coordinate}의 realification을 적용하면 전체
real dimension은
\begin{equation}
  m=2\sum_{\vk}n_o^{\vk}n_v^{\vk}.
  \label{eq:response-real-dimension}
\end{equation}
Dense real Hessian은 $m\times m$이고 float64 저장량만
\begin{equation}
  M_{\mathrm{dense}}=8m^2\ \text{bytes}
  \label{eq:dense-hessian-storage}
\end{equation}
다. Factorization과 condition-number diagnostic은 matrix copy와 workspace를 추가로
요구한다. 따라서 dense route는 작은 reference problem에서 equation-level parity를
확인하는 데 적합하지만 production scaling에는 적합하지 않다.

Dense matrix의 각 column을 얻으려면 basis vector $e_j$마다 AX를 한 번 호출해야 한다.
이는 ``$m^2$개의 equation을 푼다''는 뜻이 아니라, $m$개의 basis direction에 대한
operator action을 모아 $m^2$ element를 만든다는 뜻이다.

\subsection{Matrix-free Hessian action}

Matrix-free route는 $A$를 저장하지 않고 trial vector $u$가 들어올 때
\begin{equation}
  u\longmapsto D^{(1)}(u)
  \longmapsto \delta J[D^{(1)}]-\tfrac12\delta K[D^{(1)}]
  \longmapsto Au
  \label{eq:matrix-free-action}
\end{equation}
을 계산한다. 첫 화살표는 VO rotation을 Hermitian transition density로 만들고, 두 번째는
선택한 FFTDF/GDF/RSDF central J/K provider를 적용하며, 마지막은 AO 결과를 VO block으로
project하고 orbital-energy gap을 더한다.

따라서 Krylov iteration이 21회라면 전체 $m$-component vector를 21개만 선택해 푸는 것이
아니다. 매 iteration에서 길이 $m$인 trial vector 전체에 AX를 적용하며, 21개의 Krylov
direction이 solution을 표현하기에 충분했다는 뜻이다.

\subsection{MINRES와 energy-gap preconditioner}

Realified HF Hessian은 이상적으로 real symmetric이지만 반드시 positive definite일
필요는 없다. MINRES는 symmetric indefinite linear system에서 residual norm을 줄이는
Lanczos/Krylov method이므로 이 구조에 맞는다. $z_0=0$에서 시작한 $p$번째 subspace는
\begin{equation}
  \mathcal K_p(A,b)
  =\operatorname{span}\{b,Ab,A^2b,\ldots,A^{p-1}b\}.
  \label{eq:minres-krylov}
\end{equation}

가장 간단한 preconditioner는 independent-particle energy gap이다.
\begin{equation}
  P_{ai\vk,ai\vk}=\varepsilon_a^{\vk}-\varepsilon_i^{\vk},
  \qquad P^{-1}r\approx A^{-1}r.
  \label{eq:gap-preconditioner}
\end{equation}
이는 two-electron coupling을 버려 solution을 근사하는 것이 아니다. MINRES가 푸는
operator는 여전히 full $A$이고, $P^{-1}$은 잘못된 scaling을 줄여 subspace convergence를
빠르게 하는 coordinate transformation이다. Positive VO gap, dtype와 k-point ordering을
solver 시작 전에 검사한다.

Convergence는 solver의 internal flag만으로 승인하지 않고 실제 equation residual
\begin{equation}
  \eta_z
  =\frac{\|A_R^Tz_R+g_R\|_\infty}
         {\max(1,\|g_R\|_\infty)}
  \label{eq:scaled-z-residual}
\end{equation}
을 다시 계산한다. 필요한 경우 residual equation
\begin{equation}
  A_R^T\delta z_R=-\left(A_R^Tz_R+g_R\right),
  \qquad z_R\leftarrow z_R+\delta z_R
  \label{eq:iterative-refinement}
\end{equation}
을 같은 matrix-free solver로 풀어 iterative refinement를 수행한다.

\subsection{대칭성, complex k-point와 solver 선택}

MINRES는 realified operator가 numerical tolerance 안에서 symmetric이라는 가정에
의존한다. Reverse k-pair conjugation, $1/N_k$, occupation factor 또는 J/K transpose가
틀리면 dense test에서 $\|A_R-A_R^T\|$가 커지고 matrix-free convergence도 불안정해질 수
있다. 작은 Gamma/sampled-k case에서 dense와 matrix-free solution, gradient와 residual을
먼저 비교해야 한다.

Automatic policy는 작은 $m$에서 dense reference, 큰 $m$에서 matrix-free를 선택할 수
있다. 그러나 exact condition number를 얻으려고 큰 dense matrix를 다시 만들면 matrix-free
목적을 잃는다. Production report에서는 condition number를 ``not evaluated''로 명시하고
iteration count, residual과 minimum gap을 보존하는 편이 정직하다.

\subsection{Nuclear source batching}

Z-vector는 한 번 풀지만 $B^{Ax}$는 $3N_A$ perturbation을 갖는다. 모든 component의
AO/VO response tensor를 동시에 만들면 대략
\begin{equation}
  M_B\sim 3N_A N_kN_{\mathrm{AO}}^2\,s_{\mathrm{dtype}}
  \label{eq:nuclear-batch}
\end{equation}
로 증가한다. Component batch $\mathcal C$를 사용해
\begin{equation}
  g_{Ax}^{\mathrm{resp}}=z^TB^{Ax},\qquad (A,x)\in\mathcal C
\end{equation}
를 계산한 뒤 scalar gradient만 보존하면 결과는 변하지 않는다. Batch size는 free memory,
dtype, AO dimension과 temporary count로 정하고, OOM에는 같은 component부터 더 작은
batch로 재시도해야 한다. 이미 완료한 component를 중복 합산해서는 안 된다.

\subsection{Full derivative tensor 대신 adjoint contraction}

GDF/RSDF에서 whitened three-center factor를 $C=L^{-1}B$라 하자. Nuclear derivative는
\begin{equation}
  C^{Ax}=L^{-1}\left(B^{Ax}-L^{Ax}C\right)
  \label{eq:bounded-cderi-derivative}
\end{equation}
지만, 최종 gradient가 요구하는 것은 대개 어떤 adjoint source $W_C$와의 scalar
contraction
\begin{equation}
  \mathcal G^{Ax}=\operatorname{Re}\langle W_C,C^{Ax}\rangle
  \label{eq:reverse-contraction}
\end{equation}
뿐이다. $W_B=L^{-\dagger}W_C$를 먼저 만들면
\begin{equation}
  \mathcal G^{Ax}
  =\operatorname{Re}\langle W_B,B^{Ax}\rangle
   -\operatorname{Re}\langle W_B,L^{Ax}C\rangle.
  \label{eq:reverse-cderi-response}
\end{equation}
따라서 $(3N_A,N_{\mathrm{aux}},N_{\mathrm{AO}}^2)$ full $C^{Ax}$를 저장하지 않고 source와
metric derivative를 block별로 scalar에 누적할 수 있다. 이는 reverse-mode/adjoint
algebra이며 작은 항을 무시하는 screening approximation과 다르다.

\subsection{Signed two-density polarization}

Restricted-HF two-electron fixed-density derivative를 quadratic form
\begin{equation}
  Q^{Ax}(X)
  =\frac{1}{2N_k}\sum_{\vk}\operatorname{Re}
   \Tr\left[X^{\vk}L^{Ax}[X]^{\vk}\right]
  \label{eq:quadratic-derivative}
\end{equation}
로 쓴다. $L^{Ax}$가 density에 선형이면 reference density $D$와 Z transition density
$T$의 bilinear response는 exact polarization identity
\begin{equation}
  R_{2e}^{Z,Ax}=Q^{Ax}(D+T)-Q^{Ax}(D-T)
  \label{eq:polarization-identity}
\end{equation}
로 얻어진다. 동치인 direct form은
\begin{equation}
  R_{2e}^{Z,Ax}
  =\frac1{N_k}\sum_{\vk}\operatorname{Re}\left\{
    \Tr[D^{\vk}L^{Ax}[T]^{\vk}]
   +\Tr[T^{\vk}L^{Ax}[D]^{\vk}]
  \right\}.
  \label{eq:bilinear-response}
\end{equation}

두 개의 independent quadratic traversal 대신 signed state를 같은 SR/LR/metric traversal에
태우면 grid와 integral을 한 번 읽을 수 있다. 단, Coulomb과 exchange의 두 orientation,
auxiliary metric response, Ewald exchange correction, k-pair conjugation을 모두 보존해야
한다. 단순히 $D+T$ 한 density만 계산하면 self term이 섞여 식
\eqref{eq:bilinear-response}가 되지 않는다.

\subsection{Streaming, block size와 object lifetime}

Reciprocal source derivative가 $N_G$ vectors와 $N_{AO}^2$ AO pair를 동시에 만들면 주요
working set은 대략
\begin{equation}
  M_{\mathrm{block}}
  \sim b_GN_kN_{AO}^2s_{\mathrm{complex}}\,n_{\mathrm{temporary}}
  \label{eq:block-working-set}
\end{equation}
다. $b_G$를 줄이면 모든 $G$ vector를 그대로 처리하면서 peak memory를 낮출 수 있다.
작은 block은 kernel launch 수를 늘려 조금 느릴 수 있으므로, free-memory fraction과 hard
cap을 함께 사용하고 resolved block size/count를 report한다.

또한 마지막 consumer가 끝난 CDERI, multigrid task와 response temporary를 즉시 release하면
서로 다른 stage의 peak가 겹치지 않는다. 다음 distinction이 중요하다.
\begin{equation}
  \underbrace{\text{batching, streaming, adjoint, lifetime}}_{\text{same finite representation}}
  \neq
  \underbrace{\text{looser screening, smaller mesh, smaller basis}}_{\text{changed representation}}.
  \label{eq:algebra-vs-approx}
\end{equation}
첫 그룹은 roundoff 안에서 같은 결과를 내야 하고 parity test로 검증한다. 두 번째 그룹은
계산량과 정확도를 함께 바꾸므로 convergence study가 필요하다.

\subsection{알고리즘 참고 문헌}

MINRES의 원 논문은 C.~C. Paige와 M.~A. Saunders,
\href{https://doi.org/10.1137/0712047}{SIAM J. Numer. Anal. 12, 617--629
(1975)}다. 이 project의 matrix-free 사용은 해당 일반 알고리즘에 HF Hessian action과
orbital-gap preconditioner를 결합한 것이다. PySCF의 설계 철학과 periodic Gaussian
framework는 \href{https://doi.org/10.1002/wcms.1340}{Sun et al., WIREs Comput.
Mol. Sci. 8, e1340 (2018)}를 참고한다.

\section{계산 흐름 예제와 결과 해석}
\label{sec:worked-workflow}

이 장은 앞의 수식을 실제 restricted closed-shell Gamma 계산 순서로 연결한다. 특정
molecule의 숫자를 재현하는 input tutorial이 아니라, 어떤 object가 언제 만들어지고
어느 numerical representation을 소유하는지 추적하는 학습 예제다.

\subsection{Step 0: numerical problem을 먼저 고정한다}

예제로 source backend는 RSDF, target integrator는 \texttt{multigrid\_xc}, target
functional은 r2SCAN이라고 하자. 계산 전 다음을 고정한다.
\begin{itemize}
  \item 원자 좌표와 lattice, 단위, charge/spin,
  \item primary basis, GTH pseudopotential와 auxiliary basis,
  \item Gamma 또는 explicit k-point list와 weight,
  \item source cutoff/mesh, RSDF source/metric $\omega$와 mesh,
  \item target multigrid mesh,
  \item HF, gradient 및 Z-vector tolerance.
\end{itemize}
Source mesh와 target mesh를 모두 기록해야 ``500 Ha에서 계산했다''보다 재현 가능한
problem 정의가 된다.

\subsection{Step 1: source HF}

RSDF가 $J[D]$와 $K[D]$를 만들어
\begin{equation}
  F_{\mathrm{HF}}=h+J[D]-\tfrac12K[D]
\end{equation}
를 반복한다. 수렴 후 보존하는 핵심 object는 $D^{\mathrm{HF}}$, MO coefficient,
occupation, orbital energy와 source DF provider다. 이 단계에는 r2SCAN target
optimization이 없다. Target integrator를 uniform에서 multigrid로 바꾸어도 이 source
state가 달라지면 ownership 구현이 잘못된 것이다.

\subsection{Step 2: fixed-density target energy}

같은 $D^{\mathrm{HF}}$를 target functional에 넣는다. M1에서는
\begin{equation}
  E_{\mathrm{target}}
  =\Tr[D^{\mathrm{HF}}h]+E_J^{\mathrm{RSDF}}[D^{\mathrm{HF}}]
   +E_{\mathrm{xc}}^{\mathrm{MG}}[D^{\mathrm{HF}}]+E_{NN}.
\end{equation}
Target Fock은
\begin{equation}
  F_{\mathrm{target}}
  =h+J_{\mathrm{RSDF}}[D^{\mathrm{HF}}]
    +V_{\mathrm{xc}}^{\mathrm{MG}}[D^{\mathrm{HF}}]
\end{equation}
이다. 이 Fock로 SCF를 돌리지 않는다. Energy object는 target multigrid task와 Fock을
gradient가 재사용할 수 있도록 workspace에 보존한다.

\subsection{Step 3: nonstationary right-hand side}

Source HF orbital basis에서 target Fock의 VO block을 취한다.
\begin{equation}
  g_{ai}=2(C_v^\dagger F_{\mathrm{target}}C_o)_{ai}.
\end{equation}
$g=0$이면 target energy도 우연히 HF orbital에 stationary하므로 orbital-response term이
없다. 일반 DC-DFT에서는 $g\ne0$이고 이 값이 Z-vector equation의 right-hand side다.

\subsection{Step 4: source HF Z-vector}

Realified vector dimension $m$이 작으면 dense $A_R$를 만들어 reference solution을 얻을
수 있다. 큰 계산은 RSDF transition-density J/K를 사용하는 matrix-free AX와 MINRES를
사용한다. Report에서 다음을 확인한다.
\begin{itemize}
  \item solver가 dense인지 matrix-free인지,
  \item real response dimension과 iteration 수,
  \item minimum orbital gap과 preconditioner,
  \item final scaled residual $\eta_z$,
  \item solver 및 mean-iteration wall time.
\end{itemize}
Iteration 수가 작다는 사실만으로 정확성이 보장되지는 않는다. 실제 residual과 small-case
dense parity가 completion 근거다.

\subsection{Step 5: explicit target gradient}

Nuclear derivative를 취할 때 target density를 고정한 explicit term을 계산한다.
\begin{itemize}
  \item selected one-electron $h^{Ax}$와 GTH PP derivative,
  \item RSDF target $dJ/dR$,
  \item multigrid target $dE_{\mathrm{xc}}/dR$,
  \item target energy-weighted density와 $dS/dR$의 Pulay contraction,
  \item nuclear--nuclear derivative.
\end{itemize}
M1이므로 target $dJ/dR$는 RSDF가 소유하고 multigrid derivative는 XC만 포함한다.
M2라면 combined multigrid $dJ+dE_{\mathrm{xc}}$를 사용하고 RSDF target $dJ$를 끈다.

\subsection{Step 6: source response gradient}

각 $Ax$에 대해 one-electron, overlap 및 source RSDF two-electron nuclear source
$B^{Ax}$를 만들고 $z^TB^{Ax}$를 누적한다. Full $B$ 또는 full $C^{Ax}$를 동시에 저장할
필요는 없다. Nuclear component batching, source/metric adjoint contraction과 reciprocal
streaming을 적용한 뒤 $(N_A,3)$ scalar array만 보존한다.

\subsection{Step 7: 최종 조립}

최종 결과는 식~\eqref{eq:final-gradient-assembly}에 따라
\begin{equation}
  \mathbf g
  =\mathbf g_{\mathrm{explicit,target}}
   +\mathbf g_{\mathrm{Pulay}}
   +\mathbf g_{\mathrm{source,response}}
   +\mathbf g_{NN}
\end{equation}
로 조립한다. Public output은 force가 아니라 gradient이며 단위는 Ha/Bohr다.

\subsection{Shape와 owner 추적표}

\begin{center}
\begin{longtable}{>{\raggedright\arraybackslash}p{0.22\textwidth}
                  >{\raggedright\arraybackslash}p{0.25\textwidth}
                  >{\raggedright\arraybackslash}p{0.20\textwidth}
                  >{\raggedright\arraybackslash}p{0.22\textwidth}}
\toprule
Object & 대표 shape & Owner & 다음 consumer \\
\midrule
$D^{\mathrm{HF}}$ & $(N_k,N_{AO},N_{AO})$ & source HF & target energy, AX \\
$F_{\mathrm{target}}$ & same & target integrator + DF & VO right-hand side \\
$g_{vo}$ & k별 $(N_v,N_o)$ & common orbital layer & Z solver \\
$z_{vo}$ & k별 $(N_v,N_o)$ & source HF Hessian & nuclear response \\
Target XC gradient & $(N_A,3)$ & multigrid & final assembly \\
RSDF response & $(N_A,3)$ & RSDF adjoint & final assembly \\
Final gradient & $(N_A,3)$ & common assembly & report/FD \\
\bottomrule
\end{longtable}
\end{center}

\subsection{검증 순서}

한 번의 큰 계산이 끝났다는 사실만으로 구현이 검증되지는 않는다. 권장 순서는 다음과
같다.
\begin{enumerate}
  \item 작은 distorted geometry에서 primitive matrix/response parity를 확인한다.
  \item Gamma와 explicit one-k-point shape parity를 확인한다.
  \item sampled k-point에서 conjugation과 $1/N_k$를 확인한다.
  \item 선택한 Cartesian component의 full-energy FD를 수행한다.
  \item 전체 $3N_A$ FD 또는 representative complete small-system gate를 통과한다.
  \item source mesh, target mesh, auxiliary basis와 k-mesh를 각각 수렴시킨다.
  \item 마지막으로 larger-system timing과 memory를 측정한다.
\end{enumerate}

Analytic--FD agreement는 ``선택한 finite representation을 올바르게 미분했다''는
증거다. Mesh convergence는 ``그 finite representation이 원하는 physical limit에
충분히 가깝다''는 별도 증거다.

\subsection{Timing과 memory를 읽는 법}

End-to-end 시간을 source HF, target energy, target explicit gradient, Z-vector,
one-electron/PP, DF response와 final reporting으로 나누어 본다. GPU 계산에서 높은 CPU
사용률은 host integral generation, task construction 또는 data packing일 수 있으며 곧바로
GPU fallback을 뜻하지 않는다. 실제 owner metadata와 stage timing을 함께 본다.

Peak VRAM도 어느 stage에서 발생했는지 확인한다. OOM은 allocation failure이고
\texttt{CUDA\_ERROR\_ILLEGAL\_ADDRESS}는 실행된 kernel의 invalid access다. CUDA가
asynchronous이므로 traceback의 다음 CuPy operation이 최초 원인이 아닐 수 있다. 큰-grid
오류는 blocking diagnostic으로 first failing kernel을 찾고, 단순히 cutoff를 낮춰 숨기기
전에 physics representation과 kernel limitation을 분리한다.

\subsection{Further reading}

Periodic Gaussian implementation과 software convention은
\href{https://doi.org/10.1002/wcms.1340}{PySCF framework paper}를 출발점으로 삼는다.
RSDF의 range-separated three-center construction은 Ye와 Berkelbach,
\href{https://doi.org/10.1063/5.0046617}{J. Chem. Phys. 154, 131104 (2021)}를
참고한다. SCAN과 r2SCAN의 functional 정의는 각각
\href{https://doi.org/10.1103/PhysRevLett.115.036402}{Phys. Rev. Lett. 115,
036402 (2015)}와
\href{https://doi.org/10.1021/acs.jpclett.0c02405}{J. Phys. Chem. Lett. 11,
8208--8215 (2020)}에 있다. 이 문서의 구현-specific owner와 validation mapping은
\texttt{docs/implementation/Equation\_Code\_Map.md}를 따른다.

\appendix
\section{Tensor shape, complex convention과 equation map}
\label{sec:appendices}

\subsection{Canonical tensor shape}

아래 표는 implementation guide와 source docstring이 따라야 하는 canonical shape다.

\begin{center}
\begin{longtable}{p{0.27\textwidth}p{0.28\textwidth}p{0.34\textwidth}}
\toprule
Quantity & Shape & Convention \\
\midrule
$D^{\vk}$ & $(N_k,N_{\mathrm{AO}},N_{\mathrm{AO}})$ & occupation factor $2$ 포함 \\
$F_{\mathrm t}^{\vk}$ & $(N_k,N_{\mathrm{AO}},N_{\mathrm{AO}})$ & Hermitian \\
$F_{vo}^{\vk}$ & k별 $(N_v,N_o)$ & target Z RHS \\
$z^{\vk}$ & k별 $(N_v,N_o)$ & complex, realified solver \\
$h^{Ax,\vk}$ & $(N_A,3,N_k,N_{\mathrm{AO}},N_{\mathrm{AO}})$ & full AO matrix derivative \\
$S^{Ax,\vk}$ & same & full overlap derivative \\
$C^{ij}_{P\mu\nu}$ & pair별 $(N_{\mathrm{aux}},N_{\mathrm{AO}},N_{\mathrm{AO}})$ & canonical k-pair \\
$C^{ij,Ax}_{P\mu\nu}$ & atom/axis + CDERI shape & source+metric response \\
$g_{Ax}$ & $(N_A,3)$ & Ha/Bohr, force 아님 \\
\bottomrule
\end{longtable}
\end{center}

Gamma에서는 singleton k-axis를 내부적으로 유지한 뒤 public PySCF shape로 복원한다. 이
원칙은 Gamma와 k-point implementation이 다른 수식을 갖는 것을 막는다.

\subsection{Complex-to-real Hessian block}

Complex response $U=X+iY$에 대해 real-linear map이
\begin{equation}
  \delta F(U)=\mathcal A U+\mathcal B U^*
  \label{eq:appendix-complex-map}
\end{equation}
형태라면 realified operator는
\begin{equation}
  \begin{pmatrix}
    \operatorname{Re}\delta F\\
    \operatorname{Im}\delta F
  \end{pmatrix}
  =
  \begin{pmatrix}
    \operatorname{Re}(\mathcal A+\mathcal B) &
    -\operatorname{Im}(\mathcal A-\mathcal B)\\
    \operatorname{Im}(\mathcal A+\mathcal B) &
    \operatorname{Re}(\mathcal A-\mathcal B)
  \end{pmatrix}
  \begin{pmatrix}X\\Y\end{pmatrix}.
  \label{eq:appendix-real-hessian}
\end{equation}
Production code는 explicit four-block matrix를 항상 만들 필요는 없고 matrix-vector
product로 같은 map을 제공할 수 있다.

\subsection{Cholesky response gauge 요약}

$M=LL^\dagger$와 $W=L^{-1}M^{Ax}L^{-\dagger}$에 대해
\begin{equation}
  \Phi(W)_{ij}
  =\begin{cases}
    W_{ij},&i>j,\\
    \tfrac12W_{ii},&i=j,\\
    0,&i<j
  \end{cases}
  \label{eq:appendix-cholesky-phi}
\end{equation}
이고 $L^{Ax}=L\Phi(W)$다. Metric이 complex Hermitian이면 diagonal imaginary noise는
tolerance 안에서 제거하되 큰 anti-Hermitian residual은 오류로 처리한다.

\subsection{\texorpdfstring{$G=0$와 $q=0$}{G=0 and q=0} 체크리스트}

\begin{enumerate}
  \item Hartree $G=0$ neutral convention은 central energy와 derivative가 같아야 한다.
  \item GTH local potential constant와 ionic Ewald charge를 함께 확인한다.
  \item GDF/RSDF source $vbar$와 metric correction을 분리한다.
  \item Sampled k-point diagonal pair도 $q=0$이다.
  \item Exchange \texttt{exxdiv='ewald'} overlap derivative를 포함한다.
  \item Gamma-only 성공을 일반 $q=0$ 처리의 증거로 사용하지 않는다.
\end{enumerate}

\subsection{Equation-to-code 연결 원칙}

실제 함수 이름과 validation test는
\texttt{docs/implementation/Equation\_Code\_Map.md}에서 관리한다. 각 mapping은 equation label,
public/private 함수, input/output shape, dtype, device, unit과 검증 test를 함께 기록한다.
수식과 source가 달라지면 둘 중 하나만 수정하지 않고 같은 변경에서 map과 test를 함께
갱신한다.

핵심 연결은 다음과 같다.

\begin{center}
\begin{longtable}{ll}
\toprule
Equation & Physical implementation owner \\
\midrule
\eqref{eq:dc-energy-components} & fixed-density target energy \\
\eqref{eq:zvector-equation} & common CPU/GPU adjoint response \\
\eqref{eq:final-gradient-assembly} & backend-independent gradient assembly \\
\eqref{eq:gth-nonlocal-derivative} & GTH nonlocal projector response \\
\eqref{eq:fftdf-hartree-derivative} & FFTDF J derivative \\
\eqref{eq:gdf-cderi-derivative} & GDF source/metric response \\
\eqref{eq:mgga-tau-derivative} & SCAN/r2SCAN explicit XC derivative \\
\eqref{eq:rsdf-cderi-derivative} & RSDF independent source/metric response \\
\bottomrule
\end{longtable}
\end{center}

\section{Reference와 재현성 주의}

PySCF와 GPU4PySCF source/example은 API와 tensor convention의 primary reference다. CP2K
source/data는 GTH pseudopotential, periodic $G=0$, cutoff/grid와 force formulation을
비교하는 reference로 사용하되, GPW/auxiliary-density algorithm을 PySCF GDF/RSDF 수식과
동일시하지 않는다. 외부 source를 참고할 때 license와 attribution을 확인하고 private API
version을 result metadata에 남긴다.


\end{document}
